\documentclass[a4paper,12pt]{article}
\usepackage[utf8]{inputenc}
\usepackage[T1]{fontenc}
\usepackage[french]{babel}
\usepackage{amsmath, amssymb, amsfonts}
\usepackage{graphicx}
\usepackage{enumitem}
\usepackage{tabularx}
\usepackage{booktabs}
\usepackage{geometry}
\usepackage{tikz}
\usepackage{hyperref}
\usepackage{xcolor}
\usepackage{pdflscape}
\usepackage{fontawesome5}

\usetikzlibrary{arrows, positioning, shapes.geometric, shapes.multipart, matrix, fit, backgrounds}

% Styles for the UML class diagrams (logical model and physical/table model)
\tikzset{
  umlclass/.style={
    rectangle split, rectangle split parts=2, draw=black, thick,
    rectangle split part fill={blue!12,white},
    font=\fontsize{6.3}{7.4}\selectfont, align=left,
    rectangle split part align={center,left},
    inner sep=4pt, rounded corners=1pt, text width=3.1cm
  },
  umlclasswide/.style={
    rectangle split, rectangle split parts=2, draw=black, thick,
    rectangle split part fill={blue!12,white},
    font=\fontsize{6.3}{7.4}\selectfont, align=left,
    rectangle split part align={center,left},
    inner sep=4pt, rounded corners=1pt, text width=3.7cm
  },
  umltable/.style={
    rectangle split, rectangle split parts=2, draw=black, thick,
    rectangle split part fill={yellow!18,white},
    font=\fontsize{6.3}{7.4}\selectfont, align=left,
    rectangle split part align={center,left},
    inner sep=4pt, rounded corners=1pt, text width=3.1cm
  },
  umltablewide/.style={
    rectangle split, rectangle split parts=2, draw=black, thick,
    rectangle split part fill={yellow!18,white},
    font=\fontsize{6.3}{7.4}\selectfont, align=left,
    rectangle split part align={center,left},
    inner sep=4pt, rounded corners=1pt, text width=3.7cm
  },
  assoc/.style={draw=black!70, thick},
  mult/.style={font=\fontsize{7}{8}\selectfont, fill=white, inner sep=1.5pt}
}

\geometry{top=2cm, bottom=2cm, left=1.5cm, right=1.5cm}

% Custom commands for professional sequence diagrams with stickmen
\newcommand{\drawactor}[3]{
  % #1: name of coordinate, #2: label, #3: height of lifeline
  \begin{scope}[shift={(#1)}]
    \draw[thick, fill=blue!10] (0,1.8) circle (0.20); % head
    \draw[thick] (0,1.6) -- (0,0.9); % body
    \draw[thick] (-0.35,1.3) -- (0.35,1.3); % arms
    \draw[thick] (0,0.9) -- (-0.2,0.4); % left leg
    \draw[thick] (0,0.9) -- (0.2,0.4); % right leg
    \node[font=\bfseries, anchor=north] (#1_label) at (0, 0.15) {#2};
  \end{scope}
  \coordinate (#1_top) at (#1_label.south);
  \draw[dashed, thick, black!40] (#1_top) -- ++(0,-#3) coordinate (#1_bot);
}

\newcommand{\drawapp}[3]{
  % #1: coordinate/node name, #2: label, #3: height of lifeline
  \node[draw, fill=green!15, minimum width=3.2cm, minimum height=0.8cm, rounded corners, font=\bfseries] (#1_node) at (#1) {#2};
  \coordinate (#1_top) at (#1_node.south);
  \draw[dashed, thick, black!40] (#1_top) -- ++(0,-#3) coordinate (#1_bot);
}

% Additional lifeline styles for the "open box" (method-level) sequence diagrams.
% Internal Spring components (Controller/Service) vs. persistence vs. external systems
% get distinct colors so a reader can tell at a glance what is application code,
% what is the database, and what is a third-party system (SeaweedFS, Keycloak, SMTP...).
\newcommand{\drawcomponent}[4]{
  % #1: coordinate/node name, #2: label, #3: width, #4: height of lifeline
  \node[draw, fill=blue!12, minimum width=#3, minimum height=0.8cm, rounded corners, font=\bfseries\footnotesize, align=center] (#1_node) at (#1) {#2};
  \coordinate (#1_top) at (#1_node.south);
  \draw[dashed, thick, black!40] (#1_top) -- ++(0,-#4) coordinate (#1_bot);
}
\newcommand{\drawpersistence}[4]{
  % #1: coordinate/node name, #2: label, #3: width, #4: height of lifeline
  \node[draw, fill=yellow!25, minimum width=#3, minimum height=0.8cm, rounded corners, font=\bfseries\footnotesize, align=center] (#1_node) at (#1) {#2};
  \coordinate (#1_top) at (#1_node.south);
  \draw[dashed, thick, black!40] (#1_top) -- ++(0,-#4) coordinate (#1_bot);
}
\newcommand{\drawexternal}[4]{
  % #1: coordinate/node name, #2: label, #3: width, #4: height of lifeline
  \node[draw, fill=orange!20, minimum width=#3, minimum height=0.8cm, rounded corners, font=\bfseries\footnotesize, align=center] (#1_node) at (#1) {#2};
  \coordinate (#1_top) at (#1_node.south);
  \draw[dashed, thick, black!40] (#1_top) -- ++(0,-#4) coordinate (#1_bot);
}
% Synchronous call / return between two lifelines at a given y-offset from each _top.
% #1: y-offset (negative, e.g. -0.8cm) #2: source _top #3: target _top #4: label
\newcommand{\call}[4]{\draw[->, >=stealth, thick] ([yshift=#1]#2) -- node[above, font=\tiny, align=center] {#4} ([yshift=#1]#3);}
\newcommand{\retn}[4]{\draw[<-, >=stealth, dashed, thick] ([yshift=#1]#2) -- node[above, font=\tiny, align=center] {#4} ([yshift=#1]#3);}
% Exception propagation: red dashed return, used throughout the "Cas d'Échec" section.
\newcommand{\throwexc}[4]{\draw[<-, >=stealth, dashed, thick, red!70!black] ([yshift=#1]#2) -- node[above, font=\tiny\ttfamily, align=center, red!70!black] {#4} ([yshift=#1]#3);}

% Styles for the deployment / infrastructure diagram
\tikzset{
  depnode/.style={
    draw=black, thick, rounded corners=3pt, minimum width=2.6cm, minimum height=1.05cm,
    align=center, font=\footnotesize\bfseries, text width=2.4cm, fill=#1
  },
  depnode/.default=blue!12,
  deplink/.style={->, >=stealth, thick, black!65},
  depzone/.style={draw=black!55, thick, dashed, rounded corners=6pt, inner sep=10pt, fill=black!3}
}

\title{\textbf{Rapport de Stage \\ Système de Gestion de l'Internat AIAC \\ (ResiAIAC)}}
\author{}
\date{}

\begin{document}

\begin{titlepage}
    \centering

    \vspace*{1cm}

    \vfill

    {\LARGE \textbf{Rapport de Stage}\par}
    \vspace{0.5cm}
    {\LARGE \textbf{Système de Gestion de l'Internat AIAC}\par}
    \vspace{0.3cm}
    {\LARGE \textbf{(ResiAIAC)}\par}

    \vspace{0.8cm}

    % Horizontal separator
    \rule{0.5\textwidth}{0.1pt}

    \vspace{0.8cm}

    {\large Amine Sidki\par}
    \vspace{0.3cm}
    {\large Mohammed Yacine Daher\par}

    \vfill

    % Project description
    {\large \textit{Conception et développement d'un système distribué pour la gestion de la résidence des étudiants à l'AIAC}\par}

    \vspace{1cm}

    \vspace*{1cm}

\end{titlepage}

\tableofcontents
\newpage

\section{Étude bibliographique}

\subsection{Introduction}
Cette analyse bibliographique présente un panorama complet des technologies, frameworks, outils et concepts modernes de développement logiciel mis en œuvre dans le cadre du projet \textbf{ResiAIAC}. Ce système intègre des composants backend et frontend avancés, des architectures distribuées hautement scalables et des pratiques d'intégration et de déploiement continus (DevOps) de standard industriel.

\subsection{Technologies et Outils Backend}

\subsubsection{Spring Boot}
\textbf{Définition :} Spring Boot est un framework Java d'entreprise open-source, conçu sur l'écosystème Spring, qui élimine la complexité de configuration des applications autonomes prêtes pour la production. Il repose sur le paradigme de \textit{"convention plutôt que configuration"} (Convention over Configuration).

\textbf{Caractéristiques clés :}
\begin{itemize}[label=$\bullet$]
    \item \textbf{Auto-configuration :} Résolution et instanciation automatiques des beans requis en fonction des dépendances présentes.
    \item \textbf{Serveur Embarqué :} Intégration par défaut de serveurs d'applications (Tomcat, Jetty ou Undertow) éliminant le besoin de serveurs d'applications externes.
    \item \textbf{Gestion des Dépendances :} Utilisation des "starters" Maven/Gradle standardisant les architectures de dépendance.
\end{itemize}
\textbf{Intérêt pratique :} Permet le développement ultra-rapide de microservices extensibles et d'API REST résilientes et sécurisées.

\subsubsection{PostgreSQL}
\textbf{Définition :} Système de gestion de base de données relationnelle et objet (SGBDR) open-source de classe entreprise. Il est mondialement reconnu pour son respect rigoureux des standards SQL, sa fiabilité exceptionnelle et sa richesse fonctionnelle.

\textbf{Caractéristiques clés :}
\begin{itemize}[label=$\bullet$]
    \item \textbf{Transactions ACID :} Support complet et robuste garantissant l'intégrité absolue des données.
    \item \textbf{Modèles de Données Hybrides :} Gestion native et optimisée du format JSONB permettant d'associer la flexibilité du NoSQL à la rigueur du relationnel.
    \item \textbf{Extensibilité :} Support d'extensions avancées (PostGIS pour les données spatiales, pgvector pour l'IA).
\end{itemize}

\subsubsection{Redis}
\textbf{Définition :} Redis (Remote Dictionary Server) est une base de données en mémoire, clé-valeur, ultra-rapide. Elle est principalement exploitée comme cache applicatif, courtier de messages (Message Broker) ou stockage de sessions.

\textbf{Caractéristiques clés :}
\begin{itemize}[label=$\bullet$]
    \item Structures de données hautement optimisées (listes, hashes, sets triés).
    \item Latence inférieure à la milliseconde grâce au stockage en RAM.
    \item Persistance configurable (RDB/AOF) pour prévenir la perte accidentelle de données.
\end{itemize}

\subsubsection{SeaweedFS}
\textbf{Définition :} Système de stockage distribué d'objets, hautement scalable et efficace pour le stockage de milliards de petits fichiers, inspiré de l'architecture Haystack de Facebook.

\textbf{Caractéristiques clés :}
\begin{itemize}[label=$\bullet$]
    \item Séparation stricte de la gestion des métadonnées (Master Server) et du stockage physique des fichiers (Volume Servers).
    \item Compatibilité totale avec l'API Amazon S3.
    \item Réplication native et basculement automatique garantissant une haute disponibilité.
\end{itemize}

\textbf{Justification du choix face à MinIO :} MinIO a longtemps été la référence de facto du stockage objet auto-hébergé compatible S3. Cependant, à partir de mai 2025, l'éditeur a progressivement retiré les fonctionnalités d'administration de son édition communautaire (console de gestion, gestion des politiques et de la réplication) pour les réserver à son offre commerciale, avant de placer le dépôt open-source en mode maintenance en décembre 2025. Ce repositionnement délibéré vers l'offre entreprise «~AIStor~» a fragilisé la viabilité de MinIO comme solution open-source pérenne. SeaweedFS a donc été retenu comme l'alternative open-source la plus viable et activement maintenue, conservant une empreinte légère et une compatibilité S3 complète sans dépendre d'un éditeur en phase de repli commercial.

\subsubsection{Keycloak}
\textbf{Définition :} Solution open-source de gestion des identités et des accès (IAM), assurant l'authentification et l'autorisation via les protocoles standards OAuth2 et OpenID Connect (OIDC).

\textbf{Caractéristiques clés :}
\begin{itemize}[label=$\bullet$]
    \item \textbf{Émission de jetons JWT :} Génère des jetons signés contenant les rôles et informations de l'utilisateur authentifié, vérifiés côté backend sans appel systématique au serveur d'identité.
    \item \textbf{Gestion des \textit{clients} et des rôles :} Permet de définir des applications clientes distinctes (backend, outils de test) et une hiérarchie de rôles applicatifs.
    \item \textbf{Déploiement autonome :} Hébergé au sein du cluster Kubernetes aux côtés des services applicatifs, avec sa propre base de données dédiée.
\end{itemize}
\textbf{Intérêt pratique :} Centralise l'authentification pour un déploiement mono-tenant (une seule école), sans besoin de fédération SSO entre plusieurs fournisseurs d'identité.

\subsubsection{Spring Security \& RBAC}
\textbf{Définition :} Module de l'écosystème Spring assurant l'authentification et le contrôle d'accès aux ressources d'une application. Couplé à Keycloak, il agit en \textit{Resource Server} validant les jetons JWT entrants.

\textbf{Caractéristiques clés :}
\begin{itemize}[label=$\bullet$]
    \item \textbf{Contrôle d'accès basé sur les rôles (RBAC) :} Restriction fine des opérations selon le rôle de l'utilisateur (Étudiant, Responsable, Manager, Administrateur).
    \item \textbf{Journalisation d'audit :} Traçabilité des actions sensibles effectuées sur le système, consultable par les administrateurs.
    \item \textbf{Approvisionnement des comptes (Provisioning) :} Création conjointe du compte Keycloak et de l'entité \textit{Utilisateur} au moment de l'admission, afin d'éviter la circulation de données personnelles non chiffrées dans le jeton JWT.
\end{itemize}

\subsubsection{MapStruct}
\textbf{Définition :} Générateur de code Java, exécuté à la compilation, produisant des \textit{mappers} pour convertir automatiquement les entités JPA vers leurs Objets de Transfert de Données (DTO) et inversement.

\textbf{Caractéristiques clés :}
\begin{itemize}[label=$\bullet$]
    \item Génération de code à la compilation (pas de réflexion à l'exécution), garantissant des performances proches d'un mapping manuel.
    \item Réduction du code répétitif (\textit{boilerplate}) associé à la conversion entité/DTO sur chaque endpoint REST.
\end{itemize}

\subsubsection{JUnit 5, Mockito \& MockMvc}
\textbf{Définition :} Ensemble de bibliothèques de test formant la pile de vérification automatisée du backend : JUnit~5 comme moteur d'exécution, Mockito pour l'isolation des dépendances, AssertJ pour des assertions lisibles, et MockMvc pour tester la couche web sans démarrer de serveur réel.

\textbf{Caractéristiques clés :}
\begin{itemize}[label=$\bullet$]
    \item \textbf{Tests de service (unitaires) :} Chaque \texttt{ServiceImpl} est testé isolément, dépendances (\textit{repository}, \textit{mapper}) simulées via \texttt{@Mock} ; les appels statiques de \texttt{ResourceFetcher} sont simulés avec \texttt{mockStatic} (Mockito 5 / \texttt{mockito-inline}).
    \item \textbf{Tests de Controller (tranche web) :} Chaque Controller est testé via \texttt{@WebMvcTest} et \texttt{MockMvc}, le service associé étant remplacé par \texttt{@MockitoBean} ; seuls le routage, la (dé)sérialisation JSON et les codes de statut sont vérifiés ici, la logique métier étant déjà couverte par les tests de service.
    \item \textbf{Filtres de sécurité désactivés à ce niveau :} Les tests de Controller démarrent avec \texttt{addFilters = false}, l'authentification Keycloak/JWT étant hors périmètre de cette tranche de tests.
    \item \textbf{Couverture des champs énumérés :} Vérification explicite du comportement des champs de type \textit{enum} (ex. \texttt{EtatChambre}, \texttt{EtatDocument}), y compris leur acceptation actuelle de la valeur \texttt{NULL} — point identifié comme perfectible.
\end{itemize}
\textbf{Intérêt pratique :} Exécution automatisée à chaque contribution via GitHub Actions, avant intégration au code principal.

\subsection{Technologies et Outils Frontend}

\subsubsection{Angular}
\textbf{Définition :} Framework applicatif structuré et typé, développé par Google, utilisant TypeScript pour concevoir des applications web monopages (SPA) robustes, performantes et scalables au niveau entreprise.

\textbf{Caractéristiques clés :}
\begin{itemize}[label=$\bullet$]
    \item Architecture modulaire stricte basée sur les composants réutilisables.
    \item Injection de dépendances native pour une modularité accrue et une testabilité maximale.
    \item \textit{Two-Way Data Binding} unifiant de manière transparente le modèle et la vue.
\end{itemize}

\subsubsection{RxJS}
\textbf{Définition :} Bibliothèque réactive mettant en œuvre le pattern Observer et les flux d'observables pour orchestrer de manière asynchrone les événements et les requêtes de données.

\textbf{Caractéristiques clés :}
\begin{itemize}[label=$\bullet$]
    \item Gestion fluide des requêtes HTTP concurrentes et du temps réel (WebSockets).
    \item Large palette d'opérateurs de transformation et de filtrage (\texttt{switchMap}, \texttt{debounceTime}, \\ \texttt{combineLatest}).
\end{itemize}

\subsubsection{TailwindCSS}
\textbf{Définition :} Framework CSS orienté "utilitaires d'abord" (Utility-First), permettant de concevoir des interfaces graphiques modernes et réactives directement dans le balisage HTML.

\textbf{Caractéristiques clés :}
\begin{itemize}[label=$\bullet$]
    \item Évite la prolifération de feuilles de style complexes et redondantes.
    \item Optimisation draconienne de la taille des bundles via l'élimination automatique des classes inutilisées (PurgeCSS).
\end{itemize}

\subsection{Infrastructure, Déploiement et Concepts}

\subsubsection{Docker \& Docker Compose}
\textbf{Définition :} Docker standardise l'empaquetage des applications et de leur environnement d'exécution complet dans des conteneurs isolés, garantissant le principe \textit{"build once, run anywhere"}. Docker Compose permet la définition multi-conteneur via un unique fichier déclaratif (\texttt{compose.yml}).

\subsubsection{Kubernetes (K8s)}
\textbf{Définition :} Orchestrateur de conteneurs open-source standard de l'industrie, conçu pour automatiser le déploiement, la mise à l'échelle automatique (Auto-scaling) et la gestion des applications conteneurisées.

\subsubsection{GitHub Actions}
\textbf{Définition :} Plateforme d'automatisation CI/CD intégrée à GitHub pour exécuter des pipelines de tests automatisés, de validation de code (Linter) et de déploiement automatique sur chaque événement Git.

\subsubsection{Systèmes Distribués \& API REST}
\textbf{Définition :} Architecture distribuée structurant les services indépendamment, communiquant via le protocole HTTP sans état (Stateless) en s'appuyant sur les verbes HTTP standards.

\subsubsection{Gestion d'État (State Management)}
\textbf{Définition :} Centralisation et prévisibilité des données applicatives (par exemple via NgRx) pour assurer la cohérence des vues sur les interfaces utilisateur complexes.

\subsubsection{Mise en Cache (Caching)}
\textbf{Définition :} Stratégie d'optimisation visant à éliminer les goulots d'étranglement de la base de données en conservant temporairement en RAM les données hautement sollicitées.

\subsection{Conclusion et Références Bibliographiques}
Le choix de cette pile technologique moderne garantit la création d'une application performante, modulaire et hautement disponible.

\begin{itemize}[label=$\bullet$]
    \item \textbf{Spring Boot :} \url{https://spring.io/projects/spring-boot}
    \item \textbf{PostgreSQL :} \url{https://www.postgresql.org/docs/}
    \item \textbf{Redis :} \url{https://redis.io/docs/}
    \item \textbf{SeaweedFS :} \url{https://github.com/seaweedfs/seaweedfs}
    \item \textbf{Keycloak :} \url{https://www.keycloak.org/documentation}
    \item \textbf{Spring Security :} \url{https://docs.spring.io/spring-security/reference/}
    \item \textbf{MapStruct :} \url{https://mapstruct.org/documentation/}
    \item \textbf{Docker :} \url{https://docs.docker.com/}
    \item \textbf{Kubernetes :} \url{https://kubernetes.io/docs/}
    \item \textbf{GitHub Actions :} \url{https://docs.github.com/en/actions}
    \item \textbf{Angular :} \url{https://angular.dev/}
\end{itemize}

\newpage

\section{Étude des opportunités et Analyse de l'existant}

\subsection{Situation Actuelle et Problématique}
La gestion de la résidence universitaire de l'AIAC repose sur un écosystème hétérogène d'outils cloisonnés. L'application de gestion administrative actuelle est fermée aux étudiants, forçant ces derniers à passer par des processus manuels (cahier de réclamations physique) ou des systèmes déconnectés (formulaires BDE pour la réservation des chambres).

\subsection{Synthèse des Problèmes et Opportunités}
La fragmentation actuelle engendre plusieurs dysfonctionnements majeurs :
\begin{enumerate}
    \item \textbf{Perte de temps et d'efficacité :} La double saisie des données papier vers Excel augmente la charge administrative.
    \item \textbf{Erreurs humaines :} La transcription manuelle induit des incohérences de données.
    \item \textbf{Absence de traçabilité :} Complexité de suivi de l'historique et de l'état des réclamations.
    \item \textbf{Isolement informationnel :} Manque total de visibilité unifiée sur l'occupation de l'internat.
\end{enumerate}

\subsection{Analyse Comparative des Outils Existants}
Le tableau ci-dessous dresse l'état des lieux des solutions de contournement actuellement utilisées par les acteurs :

\begin{table}[h]
\centering
\small
\begin{tabularx}{\textwidth}{l l X l}
\toprule
\textbf{Outil Actuel} & \textbf{Acteurs} & \textbf{Objectif du Processus} & \textbf{Limites Identifiées} \\
\midrule
\textbf{Papier / Cahier} & Étudiant, Responsable & Saisie des dossiers, réclamations & Perte physique, dégradation \\
\textbf{Excel} & Manager, Responsable & Suivi des réclamations & Erreurs de formules, pas de partage \\
\textbf{Téléphone} & Logisticiens, Équipes & Prise de décision urgente & Absence d'historique écrit \\
\textbf{Email} & Collaborateurs & Envoi de pièces justificatives & Risque de non-réponse, classement \\
\textbf{WhatsApp} & Manager, Personnel & Échanges informels de terrain & Non professionnel, non structuré \\
\bottomrule
\end{tabularx}
\caption{Cartographie des outils actuels et leurs limites}
\end{table}

\newpage

\section{Conception Statistique et Modèle de Données}

Cette section pose les fondations de persistance de \textbf{ResiAIAC}~: le modèle entités-relations qui sous-tend l'ensemble des entités métier décrites plus loin dans ce document. Chaque choix de modélisation (clés composites, cardinalités, contraintes d'unicité) découle directement des règles de gestion de l'internat -- occupation des chambres, suivi des promotions, traçabilité des réclamations -- plutôt que d'un schéma générique plaqué a posteriori.

\subsection{Diagramme des Classes -- Spécifications du Modèle de Données}
Pour remplacer l'application existante et centraliser l'information, le modèle de données de \textbf{ResiAIAC} est structuré autour des entités et relations clés suivantes :

\begin{description}
    \item[\textbf{Utilisateur :}] Représente toute personne accédant au système (Étudiant, Manager, Responsable, Administrateur). Contient l'historique complet de ses réservations, de ses documents administratifs et de ses réclamations.
    \item[\textbf{Document :}] Pièce justificative déposée par l'étudiant pour son dossier d'admission (carte d'identité, attestation, etc.). Il possède un état de validation (\textit{Aucun, En attente, Validé, Invalide}) et une note sur sa validité.
    \item[\textbf{Chambre :}] Entité physique caractérisée par sa capacité, son état d'occupation et sa localisation (Étage et Bâtiment). Elle est rattachée aux réservations et réclamations.
    \item[\textbf{Reservation :}] Trace historique de l'affectation d'une chambre à un étudiant pour une période donnée.
    \item[\textbf{Reclamation :}] Ticket ouvert par un étudiant concernant un problème technique dans sa chambre (lié à une entité \textbf{Service} et à des équipements défectueux).
    \item[\textbf{UtilisateurPromotionChambre (UPC) :}] Table d'association majeure modélisant la répartition des étudiants par promotion et par chambre pour assurer un suivi historique rigoureux.
\end{description}

Le diagramme de classes UML ci-après (Figure~\ref{fig:diagramme-classes}) formalise l'ensemble de ces entités, leurs attributs ainsi que les multiplicités des relations qui les lient.

\begin{landscape}
\centering
\textbf{Figure~1 -- Diagramme de Classes UML (Modèle Logique)}\label{fig:diagramme-classes}
\vspace{0.3cm}

\resizebox{!}{0.82\textheight}{%
\begin{tikzpicture}[>=stealth]
\matrix[row sep=2.6cm, column sep=1.6cm, ampersand replacement=\&, anchor=center]{
\&\&
  \node[umlclass] (Batiment) {\textbf{Batiment}\nodepart{two} id : UUID\\nom : String\\etages : List<Etage>};
\&\&\\
\node[umlclass] (Filiere) {\textbf{Filiere}\nodepart{two} id : Long\\nom : String\\niveauMaximal : Integer\\promotions : List<Promotion>};
\&\&
  \node[umlclass] (Etage) {\textbf{Etage}\nodepart{two} id : UUID\\numero : String\\batiment : Batiment\\chambres : List<Chambre>};
\&\&\\
\node[umlclass] (Promotion) {\textbf{Promotion}\nodepart{two} id : UUID\\anneeDeDepart : Long\\anneeDeFin : Long\\niveau : Integer\\filiere : Filiere\\combinaisonsUPC : List<UPC>};
\&\&
  \node[umlclass] (Chambre) {\textbf{Chambre}\nodepart{two} id : UUID\\matricule : String\\capacite : Long\\etat : EtatChambre\\etage : Etage\\reservations : List<Reservation>\\reclamations : List<Reclamation>\\combinaisonsUPC : List<UPC>};
\&\&
  \node[umlclass] (Utilisateur) {\textbf{Utilisateur}\nodepart{two} id : UUID\\nom : String\\prenom : String\\cin : String\\adresse : String\\telephone : String\\documents : List<Document>\\reservations : List<Reservation>\\reclamations : List<Reclamation>\\combinaisonsUPC : List<UPC>\\creation : TimeStamp\\modification : TimeStamp};
\\
\node[umlclasswide] (UPC) {\textbf{UtilisateurPromotion\\Chambre}\nodepart{two} id : UPCId\\utilisateur : Utilisateur\\promotion : Promotion\\chambre : Chambre\\retard : Boolean\\note : String\\equipementsEndommages : List<EquipementUpc>};
\&\&
  \node[umlclass] (Reservation) {\textbf{Reservation}\nodepart{two} id : UUID\\utilisateur : Utilisateur\\chambre : Chambre\\etat : EtatReservation\\creation : TimeStamp};
\&\&
  \node[umlclass] (Document) {\textbf{Document}\nodepart{two} id : Long\\nomFichier : String\\nomSceau : String\\proprietaire : Utilisateur\\etat : EtatDocument\\noteSurValidite : String\\creation : TimeStamp};
\\
\node[umlclass] (Equipement) {\textbf{Equipement}\nodepart{two} id : Long\\nom : String\\reclamations : List<EquipementReclamation>\\upcs : List<EquipementUpc>};
\&
  \node[umlclass] (EquipementUpc) {\textbf{EquipementUpc}\nodepart{two} id : EquipementUpcId\\equipement : Equipement\\upc : UtilisateurPromotionChambre\\quantite : Long};
\&
  \node[umlclasswide] (Reclamation) {\textbf{Reclamation}\nodepart{two} id : UUID\\utilisateur : Utilisateur\\chambre : Chambre\\service : Service\\message : String\\equipements : List<EquipementReclamation>\\etat : EtatReclamation\\creation : TimeStamp};
\&
  \node[umlclasswide] (EquipementReclamation) {\textbf{EquipementReclamation}\nodepart{two} id : EquipementReclamationId\\equipement : Equipement\\reclamation : Reclamation\\quantite : Long};
\&
  \node[umlclass] (Service) {\textbf{Service}\nodepart{two} id : Long\\nom : String\\reclamations : List<Reclamation>};
\\
};

% Associations
\draw[assoc] (Batiment.south) -- (Etage.north) node[mult, near start, left]{1} node[mult, near end, left]{*};
\draw[assoc] (Etage.south) -- (Chambre.north) node[mult, near start, left]{1} node[mult, near end, left]{*};
\draw[assoc] (Filiere.south) -- (Promotion.north) node[mult, near start, left]{1} node[mult, near end, left]{*};
\draw[assoc] (Promotion.south) -- (UPC.north) node[mult, near start, left]{1} node[mult, near end, left]{*};
\draw[assoc] (Chambre.south west) -- (UPC.north east) node[mult, near start, above]{1} node[mult, near end, below]{*};
\draw[assoc] (Chambre.south) -- (Reservation.north) node[mult, near start, left]{1} node[mult, near end, left]{*};
\draw[assoc] (Chambre.east) -- ++(0.9,0) |- (Reclamation.north) node[mult, pos=0, above right]{1} node[mult, near end, above]{*};
\draw[assoc] (Utilisateur.south west) -- (UPC.north east) node[mult, near start, above]{1} node[mult, near end, right]{*};
\draw[assoc] (Utilisateur.south west) -- (Reservation.north east) node[mult, near start, above]{1} node[mult, near end, above]{*};
\draw[assoc] (Utilisateur.south west) -- (Reclamation.north east) node[mult, near start, right]{1} node[mult, near end, above]{*};
\draw[assoc] (Utilisateur.south) -- (Document.north) node[mult, near start, left]{1} node[mult, near end, left]{*};
\draw[assoc] (Service.north) -- ++(0,0.9) -| (Reclamation.north east) node[mult, pos=0, above right]{1} node[mult, pos=1, above left]{*};
\draw[assoc] (UPC.south east) -- (EquipementUpc.north west) node[mult, near start, above]{1} node[mult, near end, above]{*};
\draw[assoc] (Equipement.east) -- (EquipementUpc.west) node[mult, near start, above]{1} node[mult, near end, above]{*};
\draw[assoc] (Equipement.north) -- ++(0,0.9) -| (EquipementReclamation.north) node[mult, pos=0, above left]{1} node[mult, pos=1, above, yshift=1pt]{*};
\draw[assoc] (Reclamation.east) -- ++(0.4,0) |- (EquipementReclamation.north west) node[mult, pos=0, above]{1} node[mult, pos=1, above right]{*};

\end{tikzpicture}
}
\end{landscape}

\subsection{Modèle Physique de Données (Schéma de Base de Données)}
Le diagramme suivant (Figure~\ref{fig:diagramme-physique}) présente la transposition du modèle logique en un schéma physique de base de données, sous forme de diagramme de classes UML où chaque classe porte le stéréotype \guillemotleft~table~\guillemotright. Les attributs identifiants sont marqués \textbf{[PK]} (clé primaire) et les références vers d'autres tables \textbf{[FK]} (clé étrangère), conformément aux conventions de modélisation physique en UML.

\begin{landscape}
\centering
\textbf{Figure~2 -- Diagramme de Classes UML (Modèle Physique / Schéma de Base de Données)}\label{fig:diagramme-physique}
\vspace{0.3cm}

\resizebox{!}{0.82\textheight}{%
\begin{tikzpicture}[>=stealth]
\matrix[row sep=2.6cm, column sep=1.6cm, ampersand replacement=\&, anchor=center]{
\&\&
  \node[umltable] (Batiment) {\textit{\guillemotleft table\guillemotright}\\\textbf{batiment}\nodepart{two} \textbf{id : UUID [PK]}\\nom : String};
\&\&\\
\node[umltable] (Filiere) {\textit{\guillemotleft table\guillemotright}\\\textbf{filiere}\nodepart{two} \textbf{id : Long [PK]}\\nom : String\\niveauMaximal : Integer};
\&\&
  \node[umltable] (Etage) {\textit{\guillemotleft table\guillemotright}\\\textbf{etage}\nodepart{two} \textbf{id : UUID [PK]}\\batimentId : UUID [FK]\\numero : String};
\&\&\\
\node[umltable] (Promotion) {\textit{\guillemotleft table\guillemotright}\\\textbf{promotion}\nodepart{two} \textbf{id : UUID [PK]}\\filiereId : Long [FK]\\anneeDeDepart : Long\\anneeDeFin : Long\\niveau : Integer};
\&\&
  \node[umltable] (Chambre) {\textit{\guillemotleft table\guillemotright}\\\textbf{chambre}\nodepart{two} \textbf{id : UUID [PK]}\\etageId : UUID [FK]\\matricule : String\\capacite : Long\\etat : Enum};
\&\&
  \node[umltable] (Utilisateur) {\textit{\guillemotleft table\guillemotright}\\\textbf{user}\nodepart{two} \textbf{id : UUID [PK]}\\nom : String\\prenom : String\\cin : String\\adresse : String\\telephone : String\\creation : TimeStamp\\modification : TimeStamp};
\\
\node[umltablewide] (UPC) {\textit{\guillemotleft table\guillemotright}\\\textbf{user\_promotion\_chambre}\nodepart{two} \textbf{id : UUID [PK]}\\userId : UUID [FK]\\promotionId : UUID [FK]\\chambreId : UUID [FK]\\retard : Boolean\\note : String};
\&\&
  \node[umltable] (Reservation) {\textit{\guillemotleft table\guillemotright}\\\textbf{reservation}\nodepart{two} \textbf{id : UUID [PK]}\\userId : UUID [FK]\\roomId : UUID [FK]\\etat : Enum\\creation : TimeStamp};
\&\&
  \node[umltable] (Document) {\textit{\guillemotleft table\guillemotright}\\\textbf{document}\nodepart{two} \textbf{id : UUID [PK]}\\nomFichier : String\\nomSceau : String\\proprietaireId : UUID [FK]\\etat : Enum\\noteSurValidite : String\\creation : TimeStamp};
\\
\node[umltable] (Equipement) {\textit{\guillemotleft table\guillemotright}\\\textbf{equipement}\nodepart{two} \textbf{id : Long [PK]}\\nom : String};
\&
  \node[umltable] (EquipementUpc) {\textit{\guillemotleft table\guillemotright}\\\textbf{equipement\_upc}\nodepart{two} \textbf{id : UUID [PK]}\\equipementId : Long [FK]\\upcId : UUID [FK]\\quantite : Long};
\&
  \node[umltablewide] (Reclamation) {\textit{\guillemotleft table\guillemotright}\\\textbf{reclamation}\nodepart{two} \textbf{id : UUID [PK]}\\userId : UUID [FK]\\roomId : UUID [FK]\\serviceId : Long [FK]\\message : String\\etat : Enum\\creation : TimeStamp};
\&
  \node[umltablewide] (EquipementReclamation) {\textit{\guillemotleft table\guillemotright}\\\textbf{equipement\_reclamation}\nodepart{two} \textbf{id : UUID [PK]}\\idEquipement : Long [FK]\\idReclamation : UUID [FK]\\quantite : Long};
\&
  \node[umltable] (Service) {\textit{\guillemotleft table\guillemotright}\\\textbf{service}\nodepart{two} \textbf{id : Long [PK]}\\nom : String};
\\
};

% Associations (mirrors the logical model)
\draw[assoc] (Batiment.south) -- (Etage.north) node[mult, near start, left]{1} node[mult, near end, left]{*};
\draw[assoc] (Etage.south) -- (Chambre.north) node[mult, near start, left]{1} node[mult, near end, left]{*};
\draw[assoc] (Filiere.south) -- (Promotion.north) node[mult, near start, left]{1} node[mult, near end, left]{*};
\draw[assoc] (Promotion.south) -- (UPC.north) node[mult, near start, left]{1} node[mult, near end, left]{*};
\draw[assoc] (Chambre.south west) -- (UPC.north east) node[mult, near start, above]{1} node[mult, near end, below]{*};
\draw[assoc] (Chambre.south) -- (Reservation.north) node[mult, near start, left]{1} node[mult, near end, left]{*};
\draw[assoc] (Chambre.east) -- ++(0.9,0) |- (Reclamation.north) node[mult, pos=0, above right]{1} node[mult, near end, above]{*};
\draw[assoc] (Utilisateur.south west) -- (UPC.north east) node[mult, near start, above]{1} node[mult, near end, right]{*};
\draw[assoc] (Utilisateur.south west) -- (Reservation.north east) node[mult, near start, above]{1} node[mult, near end, above]{*};
\draw[assoc] (Utilisateur.south west) -- (Reclamation.north east) node[mult, near start, right]{1} node[mult, near end, above]{*};
\draw[assoc] (Utilisateur.south) -- (Document.north) node[mult, near start, left]{1} node[mult, near end, left]{*};
\draw[assoc] (Service.north) -- ++(0,0.9) -| (Reclamation.north east) node[mult, pos=0, above right]{1} node[mult, pos=1, above left]{*};
\draw[assoc] (UPC.south east) -- (EquipementUpc.north west) node[mult, near start, above]{1} node[mult, near end, above]{*};
\draw[assoc] (Equipement.east) -- (EquipementUpc.west) node[mult, near start, above]{1} node[mult, near end, above]{*};
\draw[assoc] (Equipement.north) -- ++(0,0.9) -| (EquipementReclamation.north) node[mult, pos=0, above left]{1} node[mult, pos=1, above, yshift=1pt]{*};
\draw[assoc] (Reclamation.east) -- ++(0.4,0) |- (EquipementReclamation.north west) node[mult, pos=0, above]{1} node[mult, pos=1, above right]{*};

\end{tikzpicture}
}
\end{landscape}

\newpage

\textbf{Conclusion.} Le modèle de données présenté ci-dessus est directement mappé sur les entités JPA du backend~: chaque relation, cardinalité et clé composite décrite plus haut se retrouve dans les migrations générées par Hibernate (\S9.3) et dans les DTO manipulés par la couche service, détaillée dans la section suivante.

\section{Architecture Applicative Backend}

Cette section décrit comment le modèle de données ci-dessus est exposé et manipulé côté backend~: l'organisation en couches commune aux 14 entités métier, la centralisation de la gestion d'erreurs, et les choix d'infrastructure (stockage objet, IAM) qui traversent l'ensemble de l'application.

\subsection{Organisation en Couches}
Chaque entité métier du backend (et ses opérations associées) suit rigoureusement la même organisation en couches, illustrée par la Figure~\ref{fig:architecture-couches}. Le Controller REST ne dépend que de l'\textbf{interface} du service, jamais de son implémentation concrète : ce découplage permet de faire évoluer ou remplacer la logique métier d'un service sans modifier son point d'entrée, et facilite l'écriture de tests isolés à chaque couche (voir \S1.2.8).

\begin{center}
\textbf{Figure~4 -- Organisation en Couches (Backend)}\label{fig:architecture-couches}
\vspace{0.3cm}

\begin{tikzpicture}[>=stealth]
  \node[depnode=orange!15, text width=3.2cm] (controller) at (0,7.5) {Controller\\(expose l'API REST)};
  \node[depnode=blue!12, text width=3.2cm] (serviceiface) at (0,5.5) {Service\\(interface)};
  \node[depnode=blue!22, text width=3.2cm] (serviceimpl) at (0,3.5) {ServiceImpl};
  \node[depnode=green!15, text width=2.9cm] (mapper) at (5.6,3.5) {Mapper\\(MapStruct)};
  \node[depnode=yellow!18, text width=3.2cm] (repository) at (0,1.5) {Repository\\(Spring Data JPA)};
  \node[depnode=black!8, text width=3.2cm] (db) at (0,-0.5) {Base de Données};

  \draw[deplink] (controller) -- node[right, font=\scriptsize]{appelle} (serviceiface);
  \draw[deplink, dashed] (serviceimpl) -- node[right, font=\scriptsize]{implémente} (serviceiface);
  \draw[deplink] (serviceimpl) -- node[right, font=\scriptsize]{persiste via} (repository);
  \draw[deplink] (serviceimpl) -- node[above, font=\scriptsize]{convertit via} (mapper);
  \draw[deplink] (repository) -- (db);
\end{tikzpicture}
\end{center}

\subsection{Séparation DTO / Entité}
Les entités JPA ne transitent jamais directement par l'API REST. Chaque entité expose deux objets de transfert distincts, convertis par un \textit{mapper} MapStruct dédié (voir \S1.2.7) :
\begin{itemize}[label=$\bullet$]
    \item \textbf{\texttt{XxxDto} :} représentation de sortie, retournée par les endpoints de lecture.
    \item \textbf{\texttt{XxxUpdateRequest} :} représentation d'entrée, utilisée pour la création et la mise à jour, validée avant conversion vers l'entité.
\end{itemize}
Cette séparation évite d'exposer la structure interne des entités (relations JPA, champs techniques) et permet de faire évoluer le modèle de données sans casser le contrat d'API.

\subsection{Gestion des Identifiants Composites}
Trois entités du modèle de données reposent sur une clé primaire composite, portée par une classe d'identifiant dédiée : \texttt{UtilisateurPromotionChambreId}, \texttt{EquipementUpcId} et \texttt{EquipementReclamationId}. Par convention, les méthodes \texttt{getById} et \texttt{delete} des services concernés acceptent les composants bruts de la clé (identifiants simples) et reconstituent la clé composite \textbf{au sein de la couche service}, plutôt que de laisser le Controller assembler cet objet technique. La méthode \texttt{update}, à l'inverse, accepte directement l'identifiant composite déjà construit ; cette asymétrie, actuellement assumée dans l'implémentation, est identifiée comme un point à harmoniser.

\subsection{Gestion des Erreurs}
La récupération d'une ressource par son identifiant passe systématiquement par l'utilitaire générique \texttt{ResourceFetcher.fetchResource(id, repository, resourceType)}, qui centralise la logique \textit{"récupérer ou lever une exception"} et évite sa duplication dans chaque service. Lorsque la ressource est introuvable, une \texttt{ResourceNotFoundException} (exception applicative dédiée) est levée.

Les erreurs de l'API sont uniformisées à l'échelle de l'application par un \texttt{GlobalExceptionHandler} annoté \texttt{@RestControllerAdvice}, qui centralise la traduction des exceptions applicatives en réponses HTTP structurées~: \texttt{ResourceNotFoundException} et \texttt{ResourceOwnershipMismatchException} (404), \texttt{RoomFullException} et \texttt{MethodArgumentNotValidException} (400), avec un gestionnaire générique de repli pour toute exception non anticipée (500). Cette centralisation évite la duplication de la logique de traduction d'erreur dans chaque contrôleur et garantit un format de réponse cohérent (\texttt{ErrorResponse}) à travers l'ensemble de l'API.

\textbf{Conclusion.} Cette organisation en couches et cette gestion d'erreurs centralisée constituent le socle sur lequel repose chacune des entités métier détaillées dans les diagrammes de séquence (\S7)~: tout appel entrant, quel que soit le scénario, traverse le même Controller $\rightarrow$ Service $\rightarrow$ Repository et retombe sur le même mécanisme d'erreur en cas d'échec.

\section{Architecture de Déploiement}

Cette section situe le backend applicatif décrit ci-dessus dans son environnement d'exécution cible~: comment les instances Spring Boot, Keycloak et les couches de persistance sont orchestrées, exposées et supervisées en production.
Le diagramme de déploiement suivant (Figure~\ref{fig:diagramme-deploiement}) illustre l'architecture technique cible de \textbf{ResiAIAC}, organisée en couches : couche client (SPA Angular), couche de routage/entrée (NGINX), couche applicative conteneurisée et orchestrée par Kubernetes (services Spring Boot et Keycloak pour l'IAM), couche de persistance et de stockage (PostgreSQL, Redis, SeaweedFS), couche d'observabilité (Prometheus, Grafana, la pile Logstash) et enfin la chaîne d'intégration et de déploiement continus (CI/CD).

\begin{center}
\textbf{Figure~5 -- Diagramme de Déploiement de l'Infrastructure ResiAIAC}\label{fig:diagramme-deploiement}
\vspace{0.3cm}

\resizebox{\textwidth}{!}{%
\begin{tikzpicture}[>=stealth]

  % ---------- Client layer ----------
  \node[depnode=green!15, text width=2.8cm] (client) at (0,11.2) {\faAngular\ Navigateur\\SPA Angular};

  % ---------- Edge / ingress layer ----------
  \node[depnode=orange!20, text width=2.8cm] (nginx) at (0,8.6) {\faCogs\ NGINX\\(Proxy Inverse)};

  % ---------- Kubernetes cluster ----------
  \node[depzone, fit={(-4.8,7.0) (4.8,-1.2)}, label={[font=\bfseries\small]above:Cluster Kubernetes}] (k8s) {};

  \node[depnode=blue!12, text width=3.0cm, label={[font=\bfseries]above:(x N)}] (spring) at (0,5.3) {\faServer\ Main application};
  \node[depnode=blue!12, text width=2.3cm] (keycloak) at (-2.4,2.7) {\faKey\ Keycloak\\(IAM / OAuth2)};
  \node[depnode=blue!12, text width=2.3cm] (obsvc) at (2.4,2.7) {\faHeartbeat\ Sondes\\(Metrics / Logs)};

  % ---------- Data / storage layer ----------
  \node[depnode=yellow!22, text width=2.3cm] (postgres) at (-2.9,0.0) {\faDatabase\ PostgreSQL\\(Base de Données)};
  \node[depnode=yellow!22, text width=2.3cm] (redis) at (0.1,0.0) {\faDatabase\ Redis\\(Cache)};
  \node[depnode=yellow!22, text width=2.3cm] (seaweedfs) at (3.1,0.0) {\faDatabase\ SeaweedFS\\(Stockage Objet)};

  % ---------- Observability stack ----------
  \node[depzone, fit={(6.2,7.0) (12.2,-1.2)}, label={[font=\bfseries\small]above:Observabilité}] (obs) {};
  \node[depnode=red!12, text width=2.6cm] (prometheus) at (9.2,5.3) {\faChartLine\ Prometheus\\(Métriques)};
  \node[depnode=red!12, text width=2.6cm] (grafana) at (9.2,2.7) {\faChartPie\ Grafana\\(Tableaux de Bord)};
  \node[depnode=red!12, text width=2.6cm] (logstash) at (9.2,0.0) {\faList\ Logstash\\(Agrégation de Logs)};

  % ---------- Links ----------
  \draw[deplink] (client) -- node[right, font=\scriptsize] {HTTPS} (nginx);
  \draw[deplink] (nginx) -- node[right, font=\scriptsize] {HTTP} (spring);
  \draw[deplink] (spring) -- node[left, font=\scriptsize] {Auth} (keycloak);
  \draw[deplink] (spring) -- (obsvc);
  \draw[deplink] (spring) -- (postgres);
  \draw[deplink] (spring) -- (redis);
  \draw[deplink] (spring) -- (seaweedfs);

  \draw[deplink] (obsvc) -- (prometheus);
  \draw[deplink] (prometheus) -- (grafana);
  \draw[deplink] (obsvc) -- (logstash.west);

\end{tikzpicture}
}
\end{center}

\subsection{Intégration et Déploiement Continus (CI/CD)}
La chaîne CI/CD, implémentée via GitHub Actions, applique des portes de qualité strictes à chaque pull request avant intégration~:

\begin{itemize}[label=$\bullet$]
    \item \textbf{Formatage du code :} vérifié automatiquement par \texttt{Spotless} (\texttt{mvn spotless:check}) -- toute pull request dont le code ne respecte pas le formatage standardisé échoue la vérification.
    \item \textbf{Couverture de tests :} un seuil minimal de 60\% de couverture (ratio d'instructions couvertes, mesuré par \texttt{JaCoCo}) est appliqué via \texttt{mvn verify}. Le périmètre mesuré couvre les contrôleurs et validateurs ; les entités, DTO, mappers, repositories, exceptions, configuration et \textbf{la couche d'implémentation des services (\texttt{service/impl})} en sont exclus par choix -- les interfaces de service qu'elles réalisent sont déjà comptabilisées dans le ratio mesuré, ce qui évite un double comptage de la même logique métier.
\end{itemize}

\textbf{Conclusion.} L'architecture de déploiement et la chaîne CI/CD décrites ci-dessus fournissent le cadre d'exécution et de qualité dans lequel s'inscrivent les scénarios métier modélisés dans les deux sections suivantes~: chaque appel API illustré par un diagramme de séquence transite par cette infrastructure et respecte les mêmes portes de qualité avant d'être fusionné.

\newpage

\section{Modélisation des Processus}

Cette section modélise les principaux scénarios fonctionnels de \textbf{ResiAIAC}, d'abord du point de vue de l'utilisateur sous forme de diagrammes de séquence «~boîte noire~» (seule l'interaction observable entre l'acteur et l'application est représentée), puis, pour les scénarios effectivement implémentés, sous une vue «~boîte ouverte~» détaillant le cheminement interne réel jusqu'au niveau de l'appel de méthode.
Afin de garantir une cohérence graphique et conceptuelle parfaite, l'ensemble des scénarios a été homogénéisé selon le modèle standardisé et direct : \textbf{Acteur $\rightarrow$ Application}.

\subsection{Dépôt de Documents par l'Étudiant}
Ce scénario décrit le flux d'authentification et de téléchargement sécurisé des pièces justificatives de l'étudiant.

\begin{center}
\begin{tikzpicture}
    \coordinate (etud) at (0,0);
    \coordinate (app) at (9,0);
    
    \drawactor{etud}{Étudiant}{6.5cm}
    \drawapp{app}{Application}{6.5cm}
    
    % Activation boxes
    \draw[fill=white, draw=black] ([xshift=-0.1cm, yshift=-0.5cm]app_top) rectangle ([xshift=0.1cm, yshift=-6.0cm]app_top);
    \draw[fill=white, draw=black] ([xshift=-0.1cm, yshift=-0.5cm]etud_top) rectangle ([xshift=0.1cm, yshift=-6.0cm]etud_top);
    
    % Messages
    \draw[->, >=stealth, thick] ([yshift=-0.8cm]etud_top) -- node[above, font=\small] {S'authentifie} ([yshift=-0.8cm]app_top);
    
    \draw[->, >=stealth, thick] ([yshift=-1.4cm]app_top) -- ++(1.2,0) -- ++(0,-0.6) node[right, font=\footnotesize, text width=3.2cm, align=left] {Vérifie l'état de\\validation utilisateur} -- ([yshift=-2.0cm]app_top);
    
    \draw[<-, >=stealth, dashed, thick] ([yshift=-2.6cm]etud_top) -- node[above, font=\small] {Affiche la liste des documents requis} ([yshift=-2.6cm]app_top);
    
    \draw[->, >=stealth, thick] ([yshift=-3.4cm]etud_top) -- node[above, font=\small] {Télécharge le document justificatif} ([yshift=-3.4cm]app_top);
    
    \draw[->, >=stealth, thick] ([yshift=-4.0cm]app_top) -- ++(1.2,0) -- ++(0,-0.6) node[right, font=\footnotesize, text width=3.2cm, align=left] {Stocke le fichier et\\persiste les métadonnées} -- ([yshift=-4.6cm]app_top);
    
    \draw[<-, >=stealth, dashed, thick] ([yshift=-5.2cm]etud_top) -- node[above, font=\small] {Confirme le succès du téléchargement} ([yshift=-5.2cm]app_top);
\end{tikzpicture}
\end{center}

\newpage

\subsection{Traitement et Validation de Dossier par le Responsable}
Ce diagramme détaille le flux d'analyse et de validation des pièces justificatives par le responsable d'internat.

\begin{center}
\begin{tikzpicture}
    \coordinate (resp) at (0,0);
    \coordinate (app) at (9,0);
    
    \drawactor{resp}{Responsable}{11.5cm}
    \drawapp{app}{Application}{11.5cm}
    
    % Activation boxes
    \draw[fill=white, draw=black] ([xshift=-0.1cm, yshift=-0.5cm]app_top) rectangle ([xshift=0.1cm, yshift=-11.0cm]app_top);
    \draw[fill=white, draw=black] ([xshift=-0.1cm, yshift=-0.5cm]resp_top) rectangle ([xshift=0.1cm, yshift=-11.0cm]resp_top);
    
    % Messages
    \draw[->, >=stealth, thick] ([yshift=-0.8cm]resp_top) -- node[above, font=\small] {Consulte la liste des étudiants non validés} ([yshift=-0.8cm]app_top);
    
    \draw[->, >=stealth, thick] ([yshift=-1.4cm]app_top) -- ++(1.2,0) -- ++(0,-0.6) node[right, font=\footnotesize, text width=3.2cm, align=left] {Récupère les étudiants\\"Non validé"} -- ([yshift=-2.0cm]app_top);
    
    \draw[<-, >=stealth, dashed, thick] ([yshift=-2.5cm]resp_top) -- node[above, font=\small] {Retourne la liste des étudiants} ([yshift=-2.5cm]app_top);
    
    \draw[->, >=stealth, thick] ([yshift=-3.2cm]resp_top) -- node[above, font=\small] {Choisit un étudiant} ([yshift=-3.2cm]app_top);
    
    \draw[->, >=stealth, thick] ([yshift=-3.8cm]app_top) -- ++(1.2,0) -- ++(0,-0.6) node[right, font=\footnotesize, text width=3.2cm, align=left] {Récupère les informations\\de l'étudiant choisi} -- ([yshift=-4.4cm]app_top);
    
    \draw[<-, >=stealth, dashed, thick] ([yshift=-4.9cm]resp_top) -- node[above, font=\small] {Retourne les informations de l'étudiant} ([yshift=-4.9cm]app_top);
    
    \draw[->, >=stealth, thick] ([yshift=-5.6cm]resp_top) -- node[above, font=\small] {Choisit un document} ([yshift=-5.6cm]app_top);
    
    \draw[->, >=stealth, thick] ([yshift=-6.2cm]app_top) -- ++(1.2,0) -- ++(0,-0.6) node[right, font=\footnotesize, text width=3.2cm, align=left] {Récupère le document\\en base de données} -- ([yshift=-6.8cm]app_top);
    
    \draw[<-, >=stealth, dashed, thick] ([yshift=-7.3cm]resp_top) -- node[above, font=\small] {Retourne les détails du document} ([yshift=-7.3cm]app_top);
    
    \draw[->, >=stealth, thick] ([yshift=-8.0cm]resp_top) -- node[above, font=\small] {Valide ou Rejette le document} ([yshift=-8.0cm]app_top);
    
    \draw[->, >=stealth, thick] ([yshift=-8.6cm]app_top) -- ++(1.2,0) -- ++(0,-0.6) node[right, font=\footnotesize, text width=3.2cm, align=left] {Met à jour l'état\\en base de données} -- ([yshift=-9.2cm]app_top);
    
    \draw[<-, >=stealth, dashed, thick] ([yshift=-9.8cm]resp_top) -- node[above, font=\small] {Confirme la modification de l'état} ([yshift=-9.8cm]app_top);
\end{tikzpicture}
\end{center}

\newpage

\subsection{Réservation de Chambre par l'Étudiant}
Ce flux permet à l'étudiant de consulter de manière transparente l'inventaire en temps réel et de réserver son affectation.

\begin{center}
\begin{tikzpicture}
    \coordinate (etud) at (0,0);
    \coordinate (app) at (9,0);
    
    \drawactor{etud}{Étudiant}{9.5cm}
    \drawapp{app}{Application}{9.5cm}
    
    % Activation boxes
    \draw[fill=white, draw=black] ([xshift=-0.1cm, yshift=-0.5cm]app_top) rectangle ([xshift=0.1cm, yshift=-9.0cm]app_top);
    \draw[fill=white, draw=black] ([xshift=-0.1cm, yshift=-0.5cm]etud_top) rectangle ([xshift=0.1cm, yshift=-9.0cm]etud_top);
    
    % Messages
    \draw[->, >=stealth, thick] ([yshift=-0.8cm]etud_top) -- node[above, font=\small] {Se connecte} ([yshift=-0.8cm]app_top);
    
    \draw[->, >=stealth, thick] ([yshift=-1.4cm]app_top) -- ++(1.2,0) -- ++(0,-0.6) node[right, font=\footnotesize, text width=3.2cm, align=left] {Récupère les informations\\de l'étudiant} -- ([yshift=-2.0cm]app_top);
    
    \draw[<-, >=stealth, dashed, thick] ([yshift=-2.5cm]etud_top) -- node[above, font=\small] {Retourne les informations de l'étudiant} ([yshift=-2.5cm]app_top);
    
    \draw[->, >=stealth, thick] ([yshift=-3.2cm]etud_top) -- node[above, font=\small] {Consulte les chambres ouvertes} ([yshift=-3.2cm]app_top);
    
    \draw[->, >=stealth, thick] ([yshift=-3.8cm]app_top) -- ++(1.2,0) -- ++(0,-0.6) node[right, font=\footnotesize, text width=3.2cm, align=left] {Filtre et extrait les\\chambres libres} -- ([yshift=-4.4cm]app_top);
    
    \draw[<-, >=stealth, dashed, thick] ([yshift=-4.9cm]etud_top) -- node[above, font=\small] {Retourne la liste des chambres ouvertes} ([yshift=-4.9cm]app_top);
    
    \draw[->, >=stealth, thick] ([yshift=-5.6cm]etud_top) -- node[above, font=\small] {Sélectionne et valide la chambre} ([yshift=-5.6cm]app_top);
    
    \draw[->, >=stealth, thick] ([yshift=-6.2cm]app_top) -- ++(1.2,0) -- ++(0,-0.6) node[right, font=\footnotesize, text width=3.2cm, align=left] {Enregistre l'affectation\\et bloque la chambre} -- ([yshift=-6.8cm]app_top);
    
    \draw[<-, >=stealth, dashed, thick] ([yshift=-7.4cm]etud_top) -- node[above, font=\small] {Confirme et délivre le récapitulatif} ([yshift=-7.4cm]app_top);
\end{tikzpicture}
\end{center}

\newpage

\subsection{Dépôt d'une Réclamation par l'Étudiant}
Ce scénario montre le flux de déclaration standardisé au sein du tableau de bord de l'application.

\begin{center}
\begin{tikzpicture}
    \coordinate (etud) at (0,0);
    \coordinate (app) at (9,0);
    
    \drawactor{etud}{Étudiant}{7.5cm}
    \drawapp{app}{Application}{7.5cm}
    
    % Activation boxes
    \draw[fill=white, draw=black] ([xshift=-0.1cm, yshift=-0.5cm]app_top) rectangle ([xshift=0.1cm, yshift=-7.0cm]app_top);
    \draw[fill=white, draw=black] ([xshift=-0.1cm, yshift=-0.5cm]etud_top) rectangle ([xshift=0.1cm, yshift=-7.0cm]etud_top);
    
    % Messages
    \draw[->, >=stealth, thick] ([yshift=-0.8cm]etud_top) -- node[above, font=\small] {Se connecte à son espace personnel} ([yshift=-0.8cm]app_top);
    
    \draw[->, >=stealth, thick] ([yshift=-1.4cm]app_top) -- ++(1.2,0) -- ++(0,-0.6) node[right, font=\footnotesize, text width=3.2cm, align=left] {Récupère ses données\\et sa chambre} -- ([yshift=-2.0cm]app_top);
    
    \draw[<-, >=stealth, dashed, thick] ([yshift=-2.5cm]etud_top) -- node[above, font=\small] {Retourne les données de profil} ([yshift=-2.5cm]app_top);
    
    \draw[->, >=stealth, thick] ([yshift=-3.2cm]etud_top) -- node[above, font=\small] {Saisit et envoie la réclamation} ([yshift=-3.2cm]app_top);
    
    \draw[->, >=stealth, thick] ([yshift=-3.8cm]app_top) -- ++(1.2,0) -- ++(0,-0.6) node[right, font=\footnotesize, text width=3.2cm, align=left] {Enregistre l'anomalie\\en base de données} -- ([yshift=-4.4cm]app_top);
    
    \draw[->, >=stealth, thick] ([yshift=-5.0cm]app_top) -- ++(1.2,0) -- ++(0,-0.6) node[right, font=\footnotesize, text width=3.2cm, align=left] {Notifie le manager\\de service par email} -- ([yshift=-5.6cm]app_top);
    
    \draw[<-, >=stealth, dashed, thick] ([yshift=-6.2cm]etud_top) -- node[above, font=\small] {Confirme l'enregistrement à l'écran} ([yshift=-6.2cm]app_top);
\end{tikzpicture}
\end{center}

\newpage

\subsection{Suivi de l'Avancement par l'Étudiant}
Visualisation en temps réel de l'état des tickets de maintenance ou d'assistance en cours de traitement.

\begin{center}
\begin{tikzpicture}
    \coordinate (etud) at (0,0);
    \coordinate (app) at (9,0);
    
    \drawactor{etud}{Étudiant}{9.5cm}
    \drawapp{app}{Application}{9.5cm}
    
    % Activation boxes
    \draw[fill=white, draw=black] ([xshift=-0.1cm, yshift=-0.5cm]app_top) rectangle ([xshift=0.1cm, yshift=-9.0cm]app_top);
    \draw[fill=white, draw=black] ([xshift=-0.1cm, yshift=-0.5cm]etud_top) rectangle ([xshift=0.1cm, yshift=-9.0cm]etud_top);
    
    % Messages
    \draw[->, >=stealth, thick] ([yshift=-0.8cm]etud_top) -- node[above, font=\small] {Se connecte} ([yshift=-0.8cm]app_top);
    
    \draw[->, >=stealth, thick] ([yshift=-1.4cm]app_top) -- ++(1.2,0) -- ++(0,-0.6) node[right, font=\footnotesize, text width=3.2cm, align=left] {Récupère le profil\\de l'étudiant} -- ([yshift=-2.0cm]app_top);
    
    \draw[<-, >=stealth, dashed, thick] ([yshift=-2.5cm]etud_top) -- node[above, font=\small] {Retourne les données} ([yshift=-2.5cm]app_top);
    
    \draw[->, >=stealth, thick] ([yshift=-3.2cm]etud_top) -- node[above, font=\small] {Charge l'onglet "Mes réclamations"} ([yshift=-3.2cm]app_top);
    
    \draw[->, >=stealth, thick] ([yshift=-3.8cm]app_top) -- ++(1.2,0) -- ++(0,-0.6) node[right, font=\footnotesize, text width=3.2cm, align=left] {Extrait la liste des\\tickets associés} -- ([yshift=-4.4cm]app_top);
    
    \draw[<-, >=stealth, dashed, thick] ([yshift=-4.9cm]etud_top) -- node[above, font=\small] {Affiche l'ensemble des tickets} ([yshift=-4.9cm]app_top);
    
    \draw[->, >=stealth, thick] ([yshift=-5.6cm]etud_top) -- node[above, font=\small] {Sélectionne une réclamation spécifique} ([yshift=-5.6cm]app_top);
    
    \draw[->, >=stealth, thick] ([yshift=-6.2cm]app_top) -- ++(1.2,0) -- ++(0,-0.6) node[right, font=\footnotesize, text width=3.2cm, align=left] {Récupère l'historique\\et les notes techniques} -- ([yshift=-6.8cm]app_top);
    
    \draw[<-, >=stealth, dashed, thick] ([yshift=-7.4cm]etud_top) -- node[above, font=\small] {Affiche les détails complets du ticket} ([yshift=-7.4cm]app_top);
\end{tikzpicture}
\end{center}

\newpage

\subsection{Traitement et Clôture par le Manager de Résidence}
Flux opérationnel de gestion, assignation, mise à jour et résolution des réclamations par le personnel.

\begin{center}
\begin{tikzpicture}
    \coordinate (mgr) at (0,0);
    \coordinate (app) at (9,0);
    
    \drawactor{mgr}{Manager}{11.5cm}
    \drawapp{app}{Application}{11.5cm}
    
    % Activation boxes
    \draw[fill=white, draw=black] ([xshift=-0.1cm, yshift=-0.5cm]app_top) rectangle ([xshift=0.1cm, yshift=-11.0cm]app_top);
    \draw[fill=white, draw=black] ([xshift=-0.1cm, yshift=-0.5cm]mgr_top) rectangle ([xshift=0.1cm, yshift=-11.0cm]mgr_top);
    
    % Messages
    \draw[->, >=stealth, thick] ([yshift=-0.8cm]mgr_top) -- node[above, font=\small] {Consulte l'annuaire des réclamations} ([yshift=-0.8cm]app_top);
    
    \draw[->, >=stealth, thick] ([yshift=-1.4cm]app_top) -- ++(1.2,0) -- ++(0,-0.6) node[right, font=\footnotesize, text width=3.2cm, align=left] {Filtre les tickets\\ouverts et non résolus} -- ([yshift=-2.0cm]app_top);
    
    \draw[<-, >=stealth, dashed, thick] ([yshift=-2.5cm]mgr_top) -- node[above, font=\small] {Affiche la liste des tickets en attente} ([yshift=-2.5cm]app_top);
    
    \draw[->, >=stealth, thick] ([yshift=-3.2cm]mgr_top) -- node[above, font=\small] {Prend en charge une réclamation} ([yshift=-3.2cm]app_top);
    
    \draw[->, >=stealth, thick] ([yshift=-3.8cm]app_top) -- ++(1.2,0) -- ++(0,-0.6) node[right, font=\footnotesize, text width=3.2cm, align=left] {Extrait l'ensemble des\\détails techniques} -- ([yshift=-4.4cm]app_top);
    
    \draw[<-, >=stealth, dashed, thick] ([yshift=-4.9cm]mgr_top) -- node[above, font=\small] {Affiche l'interface de traitement} ([yshift=-4.9cm]app_top);
    
    \draw[->, >=stealth, thick] ([yshift=-5.6cm]mgr_top) -- node[above, font=\small] {Valide la prise en charge} ([yshift=-5.6cm]app_top);
    
    \draw[->, >=stealth, thick] ([yshift=-6.2cm]app_top) -- ++(1.2,0) -- ++(0,-0.6) node[right, font=\footnotesize, text width=3.2cm, align=left] {Met l'état à "En\\Traitement" et notifie\\l'étudiant} -- ([yshift=-6.8cm]app_top);
    
    \draw[<-, >=stealth, dashed, thick] ([yshift=-7.3cm]mgr_top) -- node[above, font=\small] {Retourne un code de succès (HTTP 200)} ([yshift=-7.3cm]app_top);
    
    \draw[->, >=stealth, thick] ([yshift=-8.0cm]mgr_top) -- node[above, font=\small] {Déclare la résolution de l'anomalie} ([yshift=-8.0cm]app_top);
    
    \draw[->, >=stealth, thick] ([yshift=-8.6cm]app_top) -- ++(1.2,0) -- ++(0,-0.6) node[right, font=\footnotesize, text width=3.2cm, align=left] {Met l'état à "Fermé -\\Traité" en base de\\données} -- ([yshift=-9.2cm]app_top);
    
    \draw[<-, >=stealth, dashed, thick] ([yshift=-9.8cm]mgr_top) -- node[above, font=\small] {Affiche l'accusé de résolution} ([yshift=-9.8cm]app_top);
\end{tikzpicture}
\end{center}

\subsection{Processus Administratif de Libération de Chambre}
Ce diagramme de séquence décrit la déclaration de sortie d'un résident et la validation d'état de la chambre.

\begin{center}
\resizebox{\textwidth}{!}{%
\begin{tikzpicture}
    \coordinate (mgr) at (0,0);
    \coordinate (app) at (9,0);

    \drawactor{mgr}{Manager}{12.5cm}
    \drawapp{app}{Application}{12.5cm}

    % Activation boxes
    \draw[fill=white, draw=black] ([xshift=-0.1cm, yshift=-0.5cm]app_top) rectangle ([xshift=0.1cm, yshift=-12.0cm]app_top);
    \draw[fill=white, draw=black] ([xshift=-0.1cm, yshift=-0.5cm]mgr_top) rectangle ([xshift=0.1cm, yshift=-12.0cm]mgr_top);

    % Messages
    \draw[->, >=stealth, thick] ([yshift=-0.8cm]mgr_top) -- node[above, font=\small] {Consulte la liste des réclamations} ([yshift=-0.8cm]app_top);

    \draw[->, >=stealth, thick] ([yshift=-1.4cm]app_top) -- ++(1.2,0) -- ++(0,-0.6) node[right, font=\footnotesize, text width=3.2cm, align=left] {Récupère les réclamations\\non résolues} -- ([yshift=-2.0cm]app_top);

    \draw[<-, >=stealth, dashed, thick] ([yshift=-2.5cm]mgr_top) -- node[above, font=\small] {Retourne la liste des réclamations ouvertes} ([yshift=-2.5cm]app_top);

    \draw[->, >=stealth, thick] ([yshift=-3.2cm]mgr_top) -- node[above, font=\small] {Choisit une réclamation pour départ} ([yshift=-3.2cm]app_top);

    \draw[->, >=stealth, thick] ([yshift=-3.8cm]app_top) -- ++(1.2,0) -- ++(0,-0.6) node[right, font=\footnotesize, text width=3.2cm, align=left] {Récupère les informations\\complémentaires} -- ([yshift=-4.4cm]app_top);

    \draw[<-, >=stealth, dashed, thick] ([yshift=-4.9cm]mgr_top) -- node[above, font=\small] {Retourne les informations complètes sur la réclamation} ([yshift=-4.9cm]app_top);

    % Self-loop action for the Manager on the left
    \draw[->, >=stealth, thick] ([yshift=-5.6cm]mgr_top) -- ++(-1.2,0) -- ++(0,-0.6) node[left, font=\footnotesize, text width=3.2cm, align=right] {Vérifie l'état\\de la chambre} -- ([yshift=-6.2cm]mgr_top);

    \draw[->, >=stealth, thick] ([yshift=-6.8cm]mgr_top) -- node[above, font=\small, text width=12.5cm, align=center] {Complète les informations sur l'état des lieux et déclenche la génération de la lettre} ([yshift=-6.8cm]app_top);

    \draw[->, >=stealth, thick] ([yshift=-7.4cm]app_top) -- ++(1.2,0) -- ++(0,-0.6) node[right, font=\footnotesize, text width=3.2cm, align=left] {Met à jour les informations\\sur l'état des lieux.} -- ([yshift=-8.0cm]app_top);

    \draw[->, >=stealth, thick] ([yshift=-8.6cm]app_top) -- ++(1.2,0) -- ++(0,-0.6) node[right, font=\footnotesize, text width=3.2cm, align=left] {Génère une lettre de départ\\avec les nouvelles infos} -- ([yshift=-9.2cm]app_top);

    \draw[->, >=stealth, thick] ([yshift=-9.8cm]app_top) -- ++(1.2,0) -- ++(0,-0.6) node[right, font=\footnotesize, text width=3.2cm, align=left] {Met à jour l'état de\\la réclamation} -- ([yshift=-10.4cm]app_top);

    \draw[<-, >=stealth, dashed, thick] ([yshift=-11.0cm]mgr_top) -- node[above, font=\small] {Lettre de départ} ([yshift=-11.0cm]app_top);
\end{tikzpicture}%
}
\end{center}

\subsection{Gestion Applicative des Profils Utilisateurs}
Contrôle complet des habilitations des comptes et annuaires par l'administrateur du système.

\begin{center}
\begin{tikzpicture}
    \coordinate (admin) at (0,0);
    \coordinate (app) at (9,0);
    
    \drawactor{admin}{Administrateur}{7.5cm}
    \drawapp{app}{Application}{7.5cm}
    
    % Activation boxes
    \draw[fill=white, draw=black] ([xshift=-0.1cm, yshift=-0.5cm]app_top) rectangle ([xshift=0.1cm, yshift=-7.0cm]app_top);
    \draw[fill=white, draw=black] ([xshift=-0.1cm, yshift=-0.5cm]admin_top) rectangle ([xshift=0.1cm, yshift=-7.0cm]admin_top);
    
    % Messages
    \draw[->, >=stealth, thick] ([yshift=-0.8cm]admin_top) -- node[above, font=\small] {Consulte la liste complète des utilisateurs} ([yshift=-0.8cm]app_top);
    
    \draw[<-, >=stealth, dashed, thick] ([yshift=-1.5cm]admin_top) -- node[above, font=\small] {Retourne la liste complète} ([yshift=-1.5cm]app_top);
    
    \node[draw, fill=yellow!10, text width=6.2cm, align=left, font=\footnotesize, right=0.3cm of admin_top, yshift=-3.1cm] (opt) {
        \textbf{Actions Administratives autorisées :} \\
        $\bullet$ Création d'un compte utilisateur \\
        $\bullet$ Édition d'un profil existant \\
        $\bullet$ Suspension/Désactivation d'accès
    };
    
    \draw[->, >=stealth, thick] ([yshift=-4.6cm]admin_top) -- node[above, font=\small] {Soumet le formulaire de modification} ([yshift=-4.6cm]app_top);
    
    \draw[->, >=stealth, thick] ([yshift=-5.2cm]app_top) -- ++(1.2,0) -- ++(0,-0.6) node[right, font=\footnotesize, text width=3.2cm, align=left] {Exécute et valide la\\mise à jour en base} -- ([yshift=-5.8cm]app_top);
    
    \draw[<-, >=stealth, dashed, thick] ([yshift=-6.4cm]admin_top) -- node[above, font=\small] {Confirme le succès de l'opération} ([yshift=-6.4cm]app_top);
\end{tikzpicture}
\end{center}

\newpage

\subsection{Gestion Fine des Rôles (RBAC)}
Ce flux standardisé permet l'assignation et la réorganisation des rôles pour sécuriser les données et processus métiers.

\begin{center}
\begin{tikzpicture}
    \coordinate (admin) at (0,0);
    \coordinate (app) at (9,0);
    
    \drawactor{admin}{Administrateur}{5.0cm}
    \drawapp{app}{Application}{5.0cm}
    
    % Activation boxes
    \draw[fill=white, draw=black] ([xshift=-0.1cm, yshift=-0.5cm]app_top) rectangle ([xshift=0.1cm, yshift=-4.5cm]app_top);
    \draw[fill=white, draw=black] ([xshift=-0.1cm, yshift=-0.5cm]admin_top) rectangle ([xshift=0.1cm, yshift=-4.5cm]admin_top);
    
    % Messages
    \draw[->, >=stealth, thick] ([yshift=-0.8cm]admin_top) -- node[above, font=\small, text width=6cm, align=center] {Ouvre la page de gestion des rôles \& affecte un rôle} ([yshift=-0.8cm]app_top);
    
    \draw[->, >=stealth, thick] ([yshift=-1.6cm]app_top) -- ++(1.2,0) -- ++(0,-0.6) node[right, font=\footnotesize, text width=3.2cm, align=left] {Vérifie les privilèges\\et recalcule les accès} -- ([yshift=-2.2cm]app_top);
    
    \draw[->, >=stealth, thick] ([yshift=-2.6cm]app_top) -- ++(1.2,0) -- ++(0,-0.6) node[right, font=\footnotesize, text width=3.2cm, align=left] {Persiste les droits\\au sein du schéma SQL} -- ([yshift=-3.2cm]app_top);
    
    \draw[<-, >=stealth, dashed, thick] ([yshift=-3.8cm]admin_top) -- node[above, font=\small] {Affiche : "Rôle attribué avec succès"} ([yshift=-3.8cm]app_top);
\end{tikzpicture}
\end{center}

\subsection{Consultation des Journaux d'Audit Sécurisés}
Ce diagramme modélise l'extraction sécurisée des logs techniques de traçabilité des opérations.

\begin{center}
\begin{tikzpicture}
    \coordinate (admin) at (0,0);
    \coordinate (app) at (9,0);
    
    \drawactor{admin}{Administrateur}{4.5cm}
    \drawapp{app}{Application}{4.5cm}
    
    % Activation boxes
    \draw[fill=white, draw=black] ([xshift=-0.1cm, yshift=-0.5cm]app_top) rectangle ([xshift=0.1cm, yshift=-4.0cm]app_top);
    \draw[fill=white, draw=black] ([xshift=-0.1cm, yshift=-0.5cm]admin_top) rectangle ([xshift=0.1cm, yshift=-4.0cm]admin_top);
    
    % Messages
    \draw[->, >=stealth, thick] ([yshift=-0.8cm]admin_top) -- node[above, font=\small] {Consulter les journaux d'audit de l'application} ([yshift=-0.8cm]app_top);
    
    \draw[->, >=stealth, thick] ([yshift=-1.5cm]app_top) -- ++(1.2,0) -- ++(0,-0.6) node[right, font=\footnotesize, text width=3.5cm, align=left] {Exécute l'extraction des logs\\sur le stockage de sécurité} -- ([yshift=-2.1cm]app_top);
    
    \draw[<-, >=stealth, dashed, thick] ([yshift=-3.0cm]admin_top) -- node[above, font=\small] {Affiche l'interface des logs d'audit} ([yshift=-3.0cm]app_top);
\end{tikzpicture}
\end{center}

\textbf{Transition vers la vue boîte ouverte.} Les diagrammes boîte noire ci-dessus couvrent l'ensemble des scénarios fonctionnels prévus pour \textbf{ResiAIAC}, y compris certains scénarios non encore implémentés côté backend à ce stade (gestion des départs, suspension de compte, gestion des rôles, consultation des journaux d'audit) -- ils font office de spécification cible. Alors que ces diagrammes modélisaient l'interaction externe visible de l'acteur, les diagrammes qui suivent ouvrent la boîte~: chaque diagramme ci-dessous détaille, jusqu'au niveau de l'appel de méthode, le cheminement réel d'une requête à travers le Controller, la couche Service et les systèmes de persistance/stockage, tel qu'il existe effectivement dans le code source actuel. Le périmètre reprend exactement les scénarios ci-dessus qui sont implémentés côté backend~: les scénarios non implémentés ci-dessus, ainsi que la distinction artificielle entre dépôt direct et dépôt depuis l'espace étudiant (fusionnée en un seul flow, les deux scénarios blackbox correspondant au même point d'entrée self-service), en sont exclus et documentés comme tels. Convention graphique~: composants Spring internes en bleu, couche de persistance en jaune, systèmes externes en orange ; les fragments \texttt{alt}/\texttt{loop} en pointillés isolent les branches conditionnelles et les boucles.

\subsection{Dépôt de Documents par l'Étudiant (ouvert)}
Version détaillée de ce scénario : chaque appel de méthode réel entre le contrôleur, la couche service et les systèmes de persistance/stockage est représenté, jusqu'à \texttt{DocumentServiceImpl.uploadMyDocument}.

\begin{center}
\resizebox{\textwidth}{!}{%
\begin{tikzpicture}
    \coordinate (etud) at (0,0);
    \coordinate (ctrl) at (3.6,0);
    \coordinate (svc) at (7.6,0);
    \coordinate (usvc) at (11.6,0);
    \coordinate (repo) at (15.6,0);
    \coordinate (ext) at (19.2,0);

    \drawactor{etud}{Étudiant}{13.2cm}
    \drawcomponent{ctrl}{Document\\Controller}{2.8cm}{13.2cm}
    \drawcomponent{svc}{Document\\ServiceImpl}{2.8cm}{13.2cm}
    \drawcomponent{usvc}{Utilisateur\\ServiceImpl}{2.8cm}{13.2cm}
    \drawpersistence{repo}{Repository\\(JPA)}{2.6cm}{13.2cm}
    \drawexternal{ext}{SeaweedFs\\Service}{2.6cm}{13.2cm}

    \call{-0.8cm}{etud_top}{ctrl_top}{POST /document/me/upload/\{type\}\\(multipart file)}
    \call{-1.7cm}{ctrl_top}{svc_top}{uploadMyDocument(jwt, fileType, file)}
    \call{-2.4cm}{svc_top}{usvc_top}{getMyEntity(jwt)}
    \call{-3.1cm}{usvc_top}{repo_top}{UtilisateurRepository\\.findByKeycloakUser(id)}
    \retn{-3.8cm}{usvc_top}{repo_top}{Optional\textless Utilisateur\textgreater}
    \retn{-4.5cm}{svc_top}{usvc_top}{Utilisateur}
    \call{-5.3cm}{svc_top}{repo_top}{DocumentRepository\\.findFirstByNomSceauAndProprietaire(...)}
    \retn{-6.0cm}{svc_top}{repo_top}{Document existant ou null}
    \call{-6.8cm}{svc_top}{repo_top}{DocumentRepository.save(nouveauDocument)}
    \retn{-7.5cm}{svc_top}{repo_top}{Document persisté (id généré)}
    \call{-8.3cm}{svc_top}{ext_top}{uploadFile(bucket, nomFichier, file)}
    \retn{-9.0cm}{svc_top}{ext_top}{OK}

    % alt fragment: only runs if a document of the same type already existed
    \draw[dashed, rounded corners, black!55] ([xshift=-0.3cm, yshift=-9.3cm]svc_top) rectangle ([xshift=0.3cm, yshift=-11.6cm]ext_top);
    \node[anchor=north west, font=\tiny\itshape, fill=white, inner sep=1pt] at ([xshift=-0.25cm, yshift=-9.32cm]svc_top) {alt~: un document du même type existait déjà};
    \call{-10.3cm}{svc_top}{repo_top}{DocumentRepository.delete(ancienDocument)}
    \call{-11.1cm}{svc_top}{ext_top}{deleteFile(ancienBucket, ancienNomFichier)}

    \retn{-12.1cm}{ctrl_top}{svc_top}{DocumentDto}
    \retn{-12.8cm}{etud_top}{ctrl_top}{200 OK (DocumentDto)}
\end{tikzpicture}%
}
\end{center}

\emph{Note~: l'ordre réel du code écrit d'abord le nouveau document (persistance + upload SeaweedFS) avant de supprimer l'ancien -- ce choix évite qu'un étudiant se retrouve sans document du type concerné si l'upload venait à échouer en cours de route.}

\newpage


\subsection{Traitement et Validation de Dossier par le Responsable (ouvert)}
Version détaillée du Scénario A2. Le contrôle d'accès réel porte sur le rôle \texttt{MANAGER} (que \texttt{RESPONSABLE} hérite via la hiérarchie de rôles), et non sur un rôle \texttt{RESPONSABLE} dédié comme le nom du scénario pourrait le laisser penser.

\begin{center}
\resizebox{0.85\textwidth}{!}{%
\begin{tikzpicture}
    \coordinate (resp) at (0,0);
    \coordinate (ctrl) at (3.8,0);
    \coordinate (svc) at (7.8,0);
    \coordinate (repo) at (11.8,0);

    \drawactor{resp}{Responsable}{12.4cm}
    \drawcomponent{ctrl}{Document\\Controller}{2.8cm}{12.4cm}
    \drawcomponent{svc}{Document\\ServiceImpl}{2.8cm}{12.4cm}
    \drawpersistence{repo}{Repository\\(JPA)}{2.6cm}{12.4cm}

    \call{-0.8cm}{resp_top}{ctrl_top}{GET /document/\{id\}\\{[}\footnotesize hasAnyRole('MANAGER'){]}}
    \call{-1.8cm}{ctrl_top}{svc_top}{getById(id)}
    \call{-2.5cm}{svc_top}{repo_top}{ResourceFetcher.fetchResource(id, ...)}
    \retn{-3.2cm}{svc_top}{repo_top}{Document (ou ResourceNotFoundException)}

    % self-call: mapper
    \draw[->, >=stealth, thick] ([yshift=-3.9cm]svc_top) -- ++(1.0,0) -- ++(0,-0.5) node[right, font=\tiny, text width=2.6cm, align=left] {documentMapper\\.toDto(entity)} -- ([yshift=-4.4cm]svc_top);

    \retn{-5.1cm}{ctrl_top}{svc_top}{DocumentDto}
    \retn{-5.8cm}{resp_top}{ctrl_top}{200 OK (DocumentDto)}

    \call{-6.8cm}{resp_top}{ctrl_top}{PUT /document/\\\{id, etat: VALIDE\textbar INVALIDE, ...\}}
    \call{-7.8cm}{ctrl_top}{svc_top}{update(id, dto)}
    \call{-8.5cm}{svc_top}{repo_top}{ResourceFetcher.fetchResource(id, ...)}
    \retn{-9.2cm}{svc_top}{repo_top}{Document}

    \draw[->, >=stealth, thick] ([yshift=-9.9cm]svc_top) -- ++(1.0,0) -- ++(0,-0.5) node[right, font=\tiny, text width=3.4cm, align=left] {documentMapper.updateEntityFromDto(dto, entity)} -- ([yshift=-10.4cm]svc_top);

    \call{-11.3cm}{svc_top}{repo_top}{DocumentRepository.save(entity)}
    \retn{-12.0cm}{ctrl_top}{svc_top}{DocumentDto (état mis à jour)}
\end{tikzpicture}%
}
\end{center}

\newpage


\subsection{Réservation de Chambre par l'Étudiant (ouvert)}
Version détaillée du Scénario B1, jusqu'à \texttt{ReservationServiceImpl.saveMy} et la délégation vers \texttt{ChambreServiceImpl} pour la transition d'état de la chambre. Le cas d'échec (chambre déjà complète) est traité séparément dans la section \textit{Scénarios Alternatifs et Cas d'Échec}.

\begin{center}
\resizebox{\textwidth}{!}{%
\begin{tikzpicture}
    \coordinate (etud) at (0,0);
    \coordinate (ctrl) at (3.6,0);
    \coordinate (svc) at (7.6,0);
    \coordinate (usvc) at (11.4,0);
    \coordinate (csvc) at (15.2,0);
    \coordinate (repo) at (19.2,0);

    \drawactor{etud}{Étudiant}{16.0cm}
    \drawcomponent{ctrl}{Reservation\\Controller}{2.8cm}{16.0cm}
    \drawcomponent{svc}{Reservation\\ServiceImpl}{2.8cm}{16.0cm}
    \drawcomponent{usvc}{Utilisateur\\ServiceImpl}{2.8cm}{16.0cm}
    \drawcomponent{csvc}{Chambre\\ServiceImpl}{2.8cm}{16.0cm}
    \drawpersistence{repo}{Repository\\(JPA)}{2.6cm}{16.0cm}

    \call{-0.8cm}{etud_top}{ctrl_top}{POST /reservation/me\\\{chambre: id, ...\}}
    \call{-1.7cm}{ctrl_top}{svc_top}{saveMy(jwt, request)}
    \call{-2.4cm}{svc_top}{usvc_top}{getMyEntity(jwt)}
    \call{-3.1cm}{usvc_top}{repo_top}{UtilisateurRepository\\.findByKeycloakUser(id)}
    \retn{-3.8cm}{usvc_top}{repo_top}{Utilisateur}
    \retn{-4.5cm}{svc_top}{usvc_top}{Utilisateur}

    \call{-5.3cm}{svc_top}{csvc_top}{getEntityById(request.chambre())}
    \call{-6.0cm}{csvc_top}{repo_top}{ResourceFetcher.fetchResource(id, ...)}
    \retn{-6.7cm}{csvc_top}{repo_top}{Chambre}
    \retn{-7.4cm}{svc_top}{csvc_top}{Chambre}

    \node[draw, fill=yellow!15, font=\tiny, align=left, text width=3.6cm, rounded corners, anchor=west] at ([xshift=0.5cm, yshift=-7.8cm]repo_node.east) {Garde~: si \texttt{chambre.etat == OCCUPEE}, lève \texttt{RoomFullException} -- voir Cas d'Échec};

    \call{-8.6cm}{svc_top}{csvc_top}{updateEtatChambre(chambre.id, nouvelEtat)}
    \draw[->, >=stealth, thick] ([yshift=-9.3cm]csvc_top) -- ++(1.0,0) -- ++(0,-0.5) node[right, font=\tiny, text width=2.6cm, align=left] {update(id, dto)\\(délégation interne)} -- ([yshift=-9.8cm]csvc_top);
    \call{-10.6cm}{csvc_top}{repo_top}{ResourceFetcher.fetchResource(id, ...)}
    \retn{-11.3cm}{csvc_top}{repo_top}{Chambre}
    \call{-12.0cm}{csvc_top}{repo_top}{ChambreRepository.save(entity)}
    \retn{-12.7cm}{svc_top}{csvc_top}{(void)}

    \call{-13.4cm}{svc_top}{repo_top}{ReservationRepository.save(nouvelleReservation)}
    \retn{-14.1cm}{svc_top}{repo_top}{Reservation persistée}
    \retn{-14.8cm}{ctrl_top}{svc_top}{ReservationDto}
    \retn{-15.5cm}{etud_top}{ctrl_top}{200 OK (ReservationDto)}
\end{tikzpicture}%
}
\end{center}

\emph{Note~: le calcul de la nouvelle transition d'état de la chambre (LIBRE/PARTIELLEMENT\_LIBRE $\rightarrow$ PARTIELLEMENT\_LIBRE/OCCUPEE selon la capacité) est effectué avant même la persistance de la réservation, directement dans \texttt{saveMy}.}

\newpage


\subsection{Dépôt de Réclamation par l'Étudiant (ouvert)}
Les Scénarios C1 et C2 du blackbox correspondent en réalité au même unique point d'entrée self-service~: \texttt{POST /reclamation/me} $\rightarrow$ \texttt{ReclamationServiceImpl.saveMy}. Ils sont donc fusionnés en un seul diagramme ouvert.

\begin{center}
\resizebox{\textwidth}{!}{%
\begin{tikzpicture}
    \coordinate (etud) at (0,0);
    \coordinate (ctrl) at (3.6,0);
    \coordinate (svc) at (7.4,0);
    \coordinate (usvc) at (11.0,0);
    \coordinate (upcsvc) at (14.8,0);
    \coordinate (eqsvc) at (18.6,0);
    \coordinate (repo) at (22.4,0);

    \drawactor{etud}{Étudiant}{17.9cm}
    \drawcomponent{ctrl}{Reclamation\\Controller}{2.8cm}{17.9cm}
    \drawcomponent{svc}{Reclamation\\ServiceImpl}{2.8cm}{17.9cm}
    \drawcomponent{usvc}{Utilisateur\\ServiceImpl}{2.8cm}{17.9cm}
    \drawcomponent{upcsvc}{UtilPromoChambre\\ServiceImpl}{3.0cm}{17.9cm}
    \drawcomponent{eqsvc}{EquipementReclam.\\ServiceImpl}{3.0cm}{17.9cm}
    \drawpersistence{repo}{Repository\\(JPA)}{2.6cm}{17.9cm}

    \call{-0.8cm}{etud_top}{ctrl_top}{POST /reclamation/me\\\{type, description, equipements[]\}}
    \call{-1.8cm}{ctrl_top}{svc_top}{saveMy(jwt, request)}
    \call{-2.5cm}{svc_top}{usvc_top}{getMyEntity(jwt)}
    \call{-3.2cm}{usvc_top}{repo_top}{UtilisateurRepository\\.findByKeycloakUser(id)}
    \retn{-3.9cm}{usvc_top}{repo_top}{Utilisateur}
    \retn{-4.6cm}{svc_top}{usvc_top}{Utilisateur}

    \call{-5.4cm}{svc_top}{upcsvc_top}{getCurrentRoomByUser(utilisateur)}
    \call{-6.1cm}{upcsvc_top}{repo_top}{UPCRepository\\.findTopByUtilisateurOrderBy...}
    \retn{-6.8cm}{upcsvc_top}{repo_top}{UPC (ou empty $\rightarrow$ ResourceNotFoundException)}
    \retn{-7.5cm}{svc_top}{upcsvc_top}{Chambre}

    \draw[->, >=stealth, thick] ([yshift=-8.2cm]svc_top) -- ++(1.0,0) -- ++(0,-0.5) node[right, font=\tiny, text width=2.6cm, align=left] {reclamationMapper\\.toEntity(dto)} -- ([yshift=-8.7cm]svc_top);
    \call{-9.4cm}{svc_top}{repo_top}{ReclamationRepository.save(entity)}
    \retn{-10.1cm}{svc_top}{repo_top}{Reclamation persistée (id généré)}

    \draw[dashed, rounded corners, black!55] ([xshift=-0.3cm, yshift=-10.5cm]svc_top) rectangle ([xshift=0.3cm, yshift=-13.5cm]repo_top);
    \node[anchor=north west, font=\tiny\itshape, fill=white, inner sep=1pt] at ([xshift=-0.25cm, yshift=-10.52cm]svc_top) {loop~: pour chaque équipement de la requête};
    \call{-11.2cm}{svc_top}{eqsvc_top}{save(new EquipementReclamationDto(...))}
    \draw[->, >=stealth, thick] ([yshift=-11.9cm]eqsvc_top) -- ++(1.0,0) -- ++(0,-0.5) node[right, font=\tiny, text width=2.8cm, align=left] {equipementReclamationMapper.toEntity(dto)} -- ([yshift=-12.4cm]eqsvc_top);
    \call{-13.3cm}{eqsvc_top}{repo_top}{EquipementReclamationRepository\\.save(entity)}

    \retn{-14.4cm}{svc_top}{eqsvc_top}{EquipementReclamationDto}

    \draw[->, >=stealth, thick] ([yshift=-15.1cm]svc_top) -- ++(1.0,0) -- ++(0,-0.5) node[right, font=\tiny, text width=2.4cm, align=left] {reclamationMapper\\.toDto(entity)} -- ([yshift=-15.6cm]svc_top);

    \retn{-16.4cm}{ctrl_top}{svc_top}{ReclamationDto}
    \retn{-17.1cm}{etud_top}{ctrl_top}{200 OK (ReclamationDto)}
\end{tikzpicture}%
}
\end{center}

\newpage


\subsection{Suivi de l'Avancement par l'Étudiant (ouvert)}
Version détaillée du Scénario C3, couvrant les deux routes self-service de consultation~: la liste paginée (\texttt{GET /reclamation/me}) et la consultation d'une réclamation précise avec vérification de propriété (\texttt{GET /reclamation/me/\{id\}}).

\begin{center}
\resizebox{0.92\textwidth}{!}{%
\begin{tikzpicture}
    \coordinate (etud) at (0,0);
    \coordinate (ctrl) at (3.6,0);
    \coordinate (svc) at (7.4,0);
    \coordinate (usvc) at (11.0,0);
    \coordinate (repo) at (14.6,0);

    \drawactor{etud}{Étudiant}{16.5cm}
    \drawcomponent{ctrl}{Reclamation\\Controller}{2.8cm}{16.5cm}
    \drawcomponent{svc}{Reclamation\\ServiceImpl}{2.8cm}{16.5cm}
    \drawcomponent{usvc}{Utilisateur\\ServiceImpl}{2.8cm}{16.5cm}
    \drawpersistence{repo}{Repository\\(JPA)}{2.6cm}{16.5cm}

    \node[anchor=west, font=\bfseries\tiny] at ([xshift=-3.6cm, yshift=-0.3cm]etud_top) {(a) Liste paginée};
    \call{-0.8cm}{etud_top}{ctrl_top}{GET /reclamation/me?page=...}
    \call{-1.6cm}{ctrl_top}{svc_top}{getAllMy(jwt, pageable)}
    \call{-2.3cm}{svc_top}{usvc_top}{getMyEntity(jwt)}
    \call{-3.0cm}{usvc_top}{repo_top}{UtilisateurRepository\\.findByKeycloakUser(id)}
    \retn{-3.7cm}{usvc_top}{repo_top}{Utilisateur}
    \retn{-4.4cm}{svc_top}{usvc_top}{Utilisateur}
    \call{-5.2cm}{svc_top}{repo_top}{ReclamationRepository\\.findAllByUtilisateur(utilisateur, pageable)}
    \retn{-5.9cm}{svc_top}{repo_top}{Page\textless Reclamation\textgreater}
    \draw[->, >=stealth, thick] ([yshift=-6.6cm]svc_top) -- ++(1.0,0) -- ++(0,-0.5) node[right, font=\tiny, text width=2.8cm, align=left] {.map(reclamationMapper::toDto)} -- ([yshift=-7.1cm]svc_top);
    \retn{-7.9cm}{ctrl_top}{svc_top}{Page\textless ReclamationDto\textgreater}
    \retn{-8.6cm}{etud_top}{ctrl_top}{200 OK}

    \draw[dotted, black!40] ([yshift=-9.2cm]etud_top) -- ([yshift=-9.2cm]repo_top);
    \node[anchor=west, font=\bfseries\tiny] at ([xshift=-3.6cm, yshift=-9.7cm]etud_top) {(b) Consultation d'une réclamation précise};

    \call{-10.2cm}{etud_top}{ctrl_top}{GET /reclamation/me/\{id\}}
    \call{-11.0cm}{ctrl_top}{svc_top}{getMyById(jwt, id)}
    \call{-11.7cm}{svc_top}{repo_top}{ResourceFetcher.fetchResource(id, ...)}
    \retn{-12.4cm}{svc_top}{repo_top}{Reclamation}
    \call{-13.1cm}{svc_top}{usvc_top}{getMyEntity(jwt)}
    \retn{-13.8cm}{svc_top}{usvc_top}{Utilisateur}

    \node[draw, fill=yellow!15, font=\tiny, align=left, text width=3.4cm, rounded corners, anchor=west] at ([xshift=0.5cm, yshift=-14.2cm]repo_node.east) {Garde~: si \texttt{entity.utilisateur.id != utilisateur.id} $\rightarrow$ \texttt{ResourceOwnershipMismatchException} -- voir Cas d'Échec};

    \draw[->, >=stealth, thick] ([yshift=-14.8cm]svc_top) -- ++(1.0,0) -- ++(0,-0.5) node[right, font=\tiny, text width=2.6cm, align=left] {reclamationMapper\\.toDto(entity)} -- ([yshift=-15.3cm]svc_top);
    \retn{-16.0cm}{ctrl_top}{svc_top}{ReclamationDto}
    \retn{-16.7cm}{etud_top}{ctrl_top}{200 OK (ReclamationDto)}
\end{tikzpicture}%
}
\end{center}

\newpage

\subsection{Traitement et Clôture par le Manager de Résidence (ouvert)}
Version détaillée du Scénario C4, couvrant la liste des réclamations côté gestion (\texttt{GET /reclamation/}) et la mise à jour d'état (\texttt{PUT /reclamation/}).

\begin{center}
\resizebox{0.8\textwidth}{!}{%
\begin{tikzpicture}
    \coordinate (mgr) at (0,0);
    \coordinate (ctrl) at (3.8,0);
    \coordinate (svc) at (7.8,0);
    \coordinate (repo) at (11.8,0);

    \drawactor{mgr}{Manager}{14.4cm}
    \drawcomponent{ctrl}{Reclamation\\Controller}{2.8cm}{14.4cm}
    \drawcomponent{svc}{Reclamation\\ServiceImpl}{2.8cm}{14.4cm}
    \drawpersistence{repo}{Repository\\(JPA)}{2.6cm}{14.4cm}

    \node[anchor=west, font=\bfseries\tiny] at ([xshift=-3.8cm, yshift=-0.3cm]mgr_top) {(a) Liste des réclamations};
    \call{-0.8cm}{mgr_top}{ctrl_top}{GET /reclamation/?page=...\\{[}hasAnyRole('MANAGER'){]}}
    \call{-1.8cm}{ctrl_top}{svc_top}{getAll(pageable)}
    \call{-2.5cm}{svc_top}{repo_top}{ReclamationRepository.findAll(pageable)}
    \retn{-3.2cm}{svc_top}{repo_top}{Page\textless Reclamation\textgreater}
    \draw[->, >=stealth, thick] ([yshift=-3.9cm]svc_top) -- ++(1.0,0) -- ++(0,-0.5) node[right, font=\tiny, text width=2.8cm, align=left] {.map(reclamationMapper::toDto)} -- ([yshift=-4.4cm]svc_top);
    \retn{-5.2cm}{ctrl_top}{svc_top}{Page\textless ReclamationDto\textgreater}
    \retn{-5.9cm}{mgr_top}{ctrl_top}{200 OK}

    \draw[dotted, black!40] ([yshift=-6.5cm]mgr_top) -- ([yshift=-6.5cm]repo_top);
    \node[anchor=west, font=\bfseries\tiny] at ([xshift=-3.8cm, yshift=-6.9cm]mgr_top) {(b) Traitement / mise à jour d'état};

    \call{-7.9cm}{mgr_top}{ctrl_top}{PUT /reclamation/\\\{id, etat: EN\_TRAITEMENT\textbar FERME\_TRAITE\textbar ..., ...\}}
    \call{-9.1cm}{ctrl_top}{svc_top}{update(id, dto)}
    \call{-9.8cm}{svc_top}{repo_top}{ResourceFetcher.fetchResource(id, ...)}
    \retn{-10.5cm}{svc_top}{repo_top}{Reclamation}
    \draw[->, >=stealth, thick] ([yshift=-11.2cm]svc_top) -- ++(1.0,0) -- ++(0,-0.5) node[right, font=\tiny, text width=3.2cm, align=left] {reclamationMapper\\.updateEntityFromDto(dto, entity)} -- ([yshift=-11.7cm]svc_top);
    \call{-12.5cm}{svc_top}{repo_top}{ReclamationRepository.save(entity)}
    \retn{-13.2cm}{ctrl_top}{svc_top}{ReclamationDto (état mis à jour)}
    \retn{-13.9cm}{mgr_top}{ctrl_top}{200 OK (ReclamationDto)}
\end{tikzpicture}%
}
\end{center}

\textbf{Conclusion.} Ces six diagrammes couvrent l'intégralité des flux self-service et de gestion effectivement implémentés dans le backend actuel. Deux écarts avec les scénarios boîte noire d'origine ont été identifiés et corrigés au passage (fusion des deux scénarios de dépôt de réclamation en un seul, retrait du filtrage RBAC illustré pour la gestion des rôles) plutôt que reproduits tels quels. La section suivante complète cette vue par les chemins d'erreur associés à ces mêmes flux.

\newpage

\section{Scénarios Alternatifs et Cas d'Échec}

Cette section documente les chemins d'erreur réellement gérés par le code (exceptions personnalisées interceptées par \texttt{GlobalExceptionHandler}), plutôt que des cas d'échec hypothétiques. Convention~: les flèches rouges en pointillés représentent la propagation d'une exception jusqu'à sa capture.

\subsection{Échec du Dépôt de Document}
Le service \texttt{SeaweedFsService} peut lever une \texttt{BucketNotFoundException} si le bucket cible n'existe pas. Cette exception n'a \textbf{pas} de \texttt{@ExceptionHandler} dédié~: elle retombe sur le gestionnaire générique \texttt{Exception.class}, qui répond 500 avec un message générique -- un gap de granularité d'erreur identifié dans le code actuel.

\begin{center}
\resizebox{0.75\textwidth}{!}{%
\begin{tikzpicture}
    \coordinate (etud) at (0,0);
    \coordinate (ctrl) at (3.6,0);
    \coordinate (svc) at (7.4,0);
    \coordinate (ext) at (11.0,0);

    \drawactor{etud}{Étudiant}{7.0cm}
    \drawcomponent{ctrl}{Document\\Controller}{2.8cm}{7.0cm}
    \drawcomponent{svc}{Document\\ServiceImpl}{2.8cm}{7.0cm}
    \drawexternal{ext}{SeaweedFs\\Service}{2.6cm}{7.0cm}

    \call{-0.8cm}{etud_top}{ctrl_top}{POST /document/me/upload/\{type\}}
    \call{-1.6cm}{ctrl_top}{svc_top}{uploadMyDocument(...)}
    \call{-2.4cm}{svc_top}{ext_top}{uploadFile(bucket, nomFichier, file)}
    \throwexc{-3.2cm}{svc_top}{ext_top}{BucketNotFoundException}
    \throwexc{-4.4cm}{ctrl_top}{svc_top}{propagation (non catchée localement)}
    \node[draw, fill=red!8, font=\tiny, align=left, text width=4.2cm, rounded corners, anchor=west] at ([xshift=0.4cm, yshift=-5.0cm]ctrl_node.east) {GlobalExceptionHandler -- \texttt{@ExceptionHandler(Exception.class)} $\rightarrow$ 500, «~Something went wrong.~» (message générique, pas de détail sur le bucket)};
    \throwexc{-6.2cm}{etud_top}{ctrl_top}{500 Internal Server Error}
\end{tikzpicture}%
}
\end{center}

\subsection{Rejet de Document}
Il ne s'agit pas d'un échec technique mais d'une issue métier alternative légitime~: le même appel \texttt{DocumentServiceImpl.update} est utilisé, avec \texttt{etat = INVALIDE} au lieu de \texttt{VALIDE} dans le \texttt{DocumentDto} envoyé. Aucune branche de code séparée n'existe -- la validation métier (pourquoi le document est invalide) reste à la charge du responsable, hors du système.

\subsection{Chambre Complète (\texttt{RoomFullException})}
\begin{center}
\resizebox{0.85\textwidth}{!}{%
\begin{tikzpicture}
    \coordinate (etud) at (0,0);
    \coordinate (ctrl) at (3.6,0);
    \coordinate (svc) at (7.6,0);
    \coordinate (csvc) at (11.4,0);

    \drawactor{etud}{Étudiant}{7.7cm}
    \drawcomponent{ctrl}{Reservation\\Controller}{2.8cm}{7.7cm}
    \drawcomponent{svc}{Reservation\\ServiceImpl}{2.8cm}{7.7cm}
    \drawcomponent{csvc}{Chambre\\ServiceImpl}{2.8cm}{7.7cm}

    \call{-0.8cm}{etud_top}{ctrl_top}{POST /reservation/me \{chambre: id\}}
    \call{-1.6cm}{ctrl_top}{svc_top}{saveMy(jwt, request)}
    \call{-2.4cm}{svc_top}{csvc_top}{getEntityById(request.chambre())}
    \retn{-3.1cm}{svc_top}{csvc_top}{Chambre (etat = OCCUPEE)}
    \node[draw, fill=red!8, font=\tiny, align=left, text width=3.8cm, rounded corners, anchor=west] at ([xshift=0.4cm, yshift=-3.7cm]csvc_node.east) {if (chambre.getEtat() == EtatChambre.OCCUPEE) throw new RoomFullException(...)};
    \throwexc{-4.9cm}{ctrl_top}{svc_top}{RoomFullException}
    \node[draw, fill=red!8, font=\tiny, align=left, text width=4.0cm, rounded corners, anchor=west] at ([xshift=0.4cm, yshift=-5.5cm]csvc_node.east) {GlobalExceptionHandler -- \texttt{@ExceptionHandler(RoomFullException.class)} $\rightarrow$ 400};
    \throwexc{-6.7cm}{etud_top}{ctrl_top}{400 Bad Request (RoomFullException)}
\end{tikzpicture}%
}
\end{center}

\subsection{Échec de Validation d'Entrée}
La validation Bean Validation (\texttt{@Valid}) sur \texttt{ReclamationDto} s'exécute au niveau du framework Spring, \textbf{avant} que le contrôleur n'invoque le service -- le point d'échec n'est donc ni le contrôleur ni le service applicatif, mais l'infrastructure MVC elle-même.

\begin{center}
\resizebox{0.7\textwidth}{!}{%
\begin{tikzpicture}
    \coordinate (etud) at (0,0);
    \coordinate (mvc) at (3.8,0);
    \coordinate (ctrl) at (7.6,0);

    \drawactor{etud}{Étudiant}{5.0cm}
    \drawcomponent{mvc}{Spring MVC\\(Bean Validation)}{3.2cm}{5.0cm}
    \drawcomponent{ctrl}{Reclamation\\Controller}{2.8cm}{5.0cm}

    \call{-0.8cm}{etud_top}{mvc_top}{POST /reclamation/me\\(champ requis manquant)}
    \throwexc{-1.8cm}{mvc_top}{ctrl_top}{ReclamationController.saveMy(...) jamais invoqué}
    \node[draw, fill=red!8, font=\tiny, align=left, text width=4.0cm, rounded corners, anchor=west] at ([xshift=0.4cm, yshift=-2.4cm]ctrl_node.east) {GlobalExceptionHandler -- \texttt{@ExceptionHandler(\allowbreak MethodArgumentNotValidException.class)} $\rightarrow$ 400, liste des champs en erreur};
    \throwexc{-3.6cm}{etud_top}{mvc_top}{400 Bad Request (détail des champs)}
\end{tikzpicture}%
}
\end{center}

\subsection{Accès à la Réclamation d'un Autre Étudiant}
\begin{center}
\resizebox{0.85\textwidth}{!}{%
\begin{tikzpicture}
    \coordinate (etud) at (0,0);
    \coordinate (ctrl) at (3.6,0);
    \coordinate (svc) at (7.6,0);
    \coordinate (repo) at (11.6,0);

    \drawactor{etud}{Étudiant}{7.7cm}
    \drawcomponent{ctrl}{Reclamation\\Controller}{2.8cm}{7.7cm}
    \drawcomponent{svc}{Reclamation\\ServiceImpl}{2.8cm}{7.7cm}
    \drawpersistence{repo}{Repository\\(JPA)}{2.6cm}{7.7cm}

    \call{-0.8cm}{etud_top}{ctrl_top}{GET /reclamation/me/\{id\}\\(id d'un autre étudiant)}
    \call{-1.8cm}{ctrl_top}{svc_top}{getMyById(jwt, id)}
    \call{-2.5cm}{svc_top}{repo_top}{ResourceFetcher.fetchResource(id, ...)}
    \retn{-3.2cm}{svc_top}{repo_top}{Reclamation (utilisateur $\neq$ appelant)}
    \node[draw, fill=red!8, font=\tiny, align=left, text width=4.0cm, rounded corners, anchor=west] at ([xshift=0.4cm, yshift=-3.8cm]repo_node.east) {if (!entity.utilisateur.id.equals(utilisateur.id)) throw new ResourceOwnershipMismatchException(...)};
    \throwexc{-5.0cm}{ctrl_top}{svc_top}{ResourceOwnershipMismatch}
    \node[draw, fill=red!8, font=\tiny, align=left, text width=4.6cm, rounded corners, anchor=west] at ([xshift=0.4cm, yshift=-5.6cm]repo_node.east) {GlobalExceptionHandler -- \texttt{@ExceptionHandler(\allowbreak ResourceOwnershipMismatchException.class)} $\rightarrow$ 404 (et non 403 -- masque volontairement l'existence de la ressource)};
    \throwexc{-6.9cm}{etud_top}{ctrl_top}{404 Not Found}
\end{tikzpicture}%
}
\end{center}

\subsection{Transition d'État}
Vérification faite dans le code~: \texttt{ReclamationServiceImpl.update} ne comporte \textbf{aucune garde} sur la transition d'état -- n'importe quelle valeur valide de l'énumération \texttt{EtatReclamation} est acceptée quel que soit l'état courant (y compris repasser une réclamation \texttt{FERME\_TRAITE} à \texttt{EN\_ATTENTE}). Il n'y a donc pas de cas d'échec applicatif à diagrammer ici~: c'est un gap de règle métier identifié, pas une exception gérée.

\subsection{Cas Transversaux : Rejet RBAC}
Le contrôle \texttt{@PreAuthorize} s'exécute dans la chaîne de filtres Spring Security, \textbf{avant} d'atteindre la méthode du contrôleur -- commun à tous les scénarios ci-dessus lorsque le rôle du JWT ne correspond pas.

\begin{center}
\resizebox{0.7\textwidth}{!}{%
\begin{tikzpicture}
    \coordinate (usr) at (0,0);
    \coordinate (sec) at (4.6,0);
    \coordinate (ctrl) at (8.4,0);

    \drawactor{usr}{Utilisateur}{4.6cm}
    \drawcomponent{sec}{Spring Security\\(@PreAuthorize)}{3.2cm}{4.6cm}
    \drawcomponent{ctrl}{Controller}{2.8cm}{4.6cm}

    \call{-0.8cm}{usr_top}{sec_top}{Requête avec JWT (rôle ETUDIANT insuffisant)}
    \throwexc{-1.8cm}{sec_top}{ctrl_top}{AccessDeniedException (méthode jamais invoquée)}
    \throwexc{-2.9cm}{usr_top}{sec_top}{403 Forbidden}
\end{tikzpicture}%
}
\end{center}

\textbf{Conclusion.} L'ensemble des chemins d'erreur documentés ci-dessus correspond à des exceptions et gardes effectivement présentes dans le code~: aucun cas hypothétique n'a été ajouté, et à l'inverse, l'absence de garde constatée pour C4 a été documentée comme telle plutôt que masquée. Cette rigueur clôt la modélisation fonctionnelle du backend~; la section suivante situe ce travail dans le cycle de vie global du projet.

\newpage

\section{Perspectives et Périmètre du Projet}

L'état actuel de \textbf{ResiAIAC} correspond à l'achèvement de sa \textbf{Phase 1 -- Backend \& Infrastructure}~: un choix délibéré de concentrer l'effort sur la construction d'une architecture backend distribuée robuste (persistance, sécurité RBAC via Keycloak, stockage objet, gestion d'erreurs centralisée, pipeline CI/CD avec portes de qualité) avant d'étendre le périmètre applicatif. Les axes suivants sont scopés pour les phases ultérieures du cycle de vie du projet.

\subsection{Frontend}
Le socle backend étant désormais consolidé, le développement du frontend (SPA Angular, conformément à l'étude bibliographique) constitue la prochaine phase immédiate du projet.

\subsection{Contrat d'API et Documentation}
La génération dynamique d'une documentation OpenAPI/Swagger est planifiée comme amélioration future, une fois les endpoints de l'API stabilisés. Ce séquencement délibéré évite une dérive de documentation (\textit{documentation drift}) pendant la phase actuelle de prototypage rapide, où les contrats d'API sont encore susceptibles d'évoluer.

\subsection{Versionnement de Base de Données}
La gestion du schéma est actuellement assurée par JPA/Hibernate (\texttt{ddl-auto: update}), un choix privilégiant la vélocité de développement durant cette phase. L'intégration d'un outil de migration versionnée tel que Flyway ou Liquibase est prévue pour la phase de mise en production, afin de garantir un versionnement immuable et traçable du schéma de base de données.

\subsection{Stratégie de Mise en Cache (Redis)}
Redis, déjà retenu lors de l'étude bibliographique et présent dans l'architecture de déploiement cible, n'est pas encore exploité applicativement (aucune annotation \texttt{@Cacheable} n'est présente dans le code actuel). La stratégie de mise en cache ciblera, dans une phase ultérieure, les données à faible volatilité (services, équipements) ainsi que les requêtes de lecture à haute fréquence (résolution utilisateur $\rightarrow$ Keycloak, promotion/chambre courante d'un étudiant). La séparation DTO/Entité et la centralisation des accès via la couche service ont été conçues pour que l'ajout d'un cache sur un endpoint futur reste une modification localisée, à faible friction.

\newpage

\end{document}
